\documentclass[8pt]{beamer}
\usetheme{Madrid}
\usecolortheme{default}
\usepackage{amsmath, amssymb, graphicx, array, multirow, booktabs, xcolor, tikz}
\usetikzlibrary{shapes,arrows,positioning,fit,backgrounds}

% Compact font settings
\setbeamerfont{title}{size=\Large}
\setbeamerfont{subtitle}{size=\small}
\setbeamerfont{author}{size=\scriptsize}
\setbeamerfont{date}{size=\scriptsize}
\setbeamerfont{frametitle}{size=\large}
\setbeamerfont{framesubtitle}{size=\normalsize}
\setbeamerfont{enumerate item}{size=\footnotesize}
\setbeamerfont{itemize item}{size=\footnotesize}
\setbeamerfont{block title}{size=\footnotesize}

% Compact spacing
\setlength{\itemsep}{0.08ex}
\setlength{\parskip}{0.08ex}
\setlength{\leftmargini}{0.3cm}
\setbeamersize{text margin left=3mm, text margin right=3mm}
\addtobeamertemplate{frame begin}{\vspace{-0.4cm}}{}

% Colors
\definecolor{myblue}{RGB}{0,0,255}
\definecolor{myred}{RGB}{255,0,0}
\definecolor{mygreen}{RGB}{0,128,0}
\definecolor{myorange}{RGB}{255,165,0}
\definecolor{mypurple}{RGB}{128,0,128}
\definecolor{mygray}{RGB}{128,128,128}
\definecolor{mygold}{RGB}{255,215,0}

\title{Responsible \& Ethical AI}
\subtitle{Complete Tutorial: Ethical Frameworks, Principles, Bias, Governance, Regulation}
\author{GARP RAI Exam Preparation}
\date{}

\begin{document}
	
	\begin{frame}
		\titlepage
	\end{frame}
	
	% ============================================================
	% SECTION 1: INTRODUCTION TO RESPONSIBLE AI - SLIDES 1-5
	% ============================================================
	
	\begin{frame}{Slide 1: What is Responsible AI?}
		\fontsize{6.5}{7.5}\selectfont
		\textbf{Definition of Responsible AI:}
		\begin{itemize}
			\item Responsible AI refers to the ethical and responsible development, deployment, and use of artificial intelligence systems
			\item Aims to ensure that AI benefits society holistically while mitigating potential risks
			\item Involves aligning AI systems with ethical standards, industry standards, regulation, legislation, and global principles
			\item Ensures AI operates in a fair, transparent, and accountable manner
		\end{itemize}
		
		\vspace{0.1cm}
		\textbf{Key Components:}
		\begin{enumerate}
			\item \textbf{Ethical Standards:} Aligning with moral principles and values
			\item \textbf{Industry Standards:} Following best practices and guidelines
			\item \textbf{Regulation and Legislation:} Complying with laws and regulations
			\item \textbf{Global Principles:} Adhering to international norms and agreements
		\end{enumerate}
		
		\textcolor{myblue}{\textbf{Key Insight:}} Responsible AI is about \textcolor{myblue}{ethical development and deployment} of AI systems!
	\end{frame}
	
	\begin{frame}{Slide 2: Purpose of Ethical Frameworks}
		\fontsize{6.5}{7.5}\selectfont
		\textbf{Why Ethical Frameworks Matter in AI:}
		\begin{itemize}
			\item Ethical frameworks help us decide how to design AI and how to deploy it
			\item They help minimize risks, harms, and wrongs
			\item They shape AI by helping us define principles to abide by
			\item They help identify unacceptable biases
			\item They assess what kind of transparency is desirable
			\item They identify stakeholders and their interests
			\item They promote accountability
			\item They make justifiable decisions about design and implementation
		\end{itemize}
		
		\vspace{0.1cm}
		\textbf{Key Benefit:}
		Different ethical frameworks help us make justifiable decisions in pluralistic societies where people hold diverse values.
		
		\textcolor{myblue}{\textbf{Key Insight:}} Ethical frameworks \textcolor{myblue}{guide decision-making} to minimize risks and harms!
	\end{frame}
	
	\begin{frame}{Slide 3: Practical Ethics Defined}
		\fontsize{6.5}{7.5}\selectfont
		\textbf{What is Practical Ethics?}
		\begin{itemize}
			\item The application of ethical theory and the development of good practices to solve real-world moral dilemmas
			\item Provides concrete guidance for moral decision making and problem solving
			\item Goes beyond theoretical philosophy to offer actionable advice
		\end{itemize}
		
		\vspace{0.1cm}
		\textbf{Key Characteristics:}
		\begin{enumerate}
			\item \textbf{Applied Focus:}
			\begin{itemize}
				\item Addresses actual ethical dilemmas
				\item Offers recommendations, not just theories
				\item Domains: medicine, law, politics, business, AI
			\end{itemize}
			\item \textbf{Multidisciplinary:}
			\begin{itemize}
				\item Draws on moral philosophy
				\item Social sciences (psychology, sociology)
				\item Domain expertise, law, computer science
			\end{itemize}
			\item \textbf{Context-Sensitive:}
			\begin{itemize}
				\item Situational nuances matter
				\item One-size-fits-all approaches don't work
			\end{itemize}
			\item \textbf{Pluralistic:}
			\begin{itemize}
				\item Blends different ethical frameworks
				\item Aims for justifiable decisions
			\end{itemize}
		\end{enumerate}
		
		\textcolor{myblue}{\textbf{Key Insight:}} Practical ethics = \textcolor{myblue}{applying theory to real-world dilemmas}!
	\end{frame}
	
	\begin{frame}{Slide 4: Why Firms Consider Ethics Frameworks}
		\fontsize{6.5}{7.5}\selectfont
		\textbf{Business Benefits of Practical Ethics:}
		
		\vspace{0.1cm}
		\textbf{Building Trust and Reputation:}
		\begin{itemize}
			\item Customers, employees, and the public trust ethical companies
			\item Improves brand image
			\item Creates competitive advantage
		\end{itemize}
		
		\vspace{0.1cm}
		\textbf{Avoiding Scandals:}
		\begin{itemize}
			\item Unethical behavior damages reputation
			\item Fraud, bias, privacy violations can be catastrophic
			\item Ethics helps prevent major scandals
		\end{itemize}
		
		\vspace{0.1cm}
		\textbf{Attracting and Retaining Talent:}
		\begin{itemize}
			\item Workers increasingly care about ethical practices
			\item Ethics frameworks help recruit and retain talent
			\item Values-aligned workforce
		\end{itemize}
		
		\textcolor{myblue}{\textbf{Key Insight:}} Ethics makes \textcolor{myblue}{good business sense} with multiple benefits!
	\end{frame}
	
	\begin{frame}{Slide 5: Ethics and Law Relationship}
		\fontsize{6.5}{7.5}\selectfont
		\textbf{How Ethics and Law Relate:}
		
		\vspace{0.1cm}
		\textbf{Ethics as Complement to Law:}
		\begin{itemize}
			\item Ethics helps ground, inform, and shape law
			\item Societies regulate behavior based on moral acceptability
			\item Ethics helps distinguish between just and unjust laws
		\end{itemize}
		
		\vspace{0.1cm}
		\textbf{Laws are Narrow in Scope:}
		\begin{itemize}
			\item Establish minimal behavioral requirements
			\item Enable orderly interactions within basic fairness
			\item Provide a framework for society
		\end{itemize}
		
		\vspace{0.1cm}
		\textbf{Ethics Goes Beyond Law:}
		\begin{itemize}
			\item Identifies moral issues
			\item Reflects on the kind of society we want
			\item Makes recommendations for a good life
			\item More ambitious than law
		\end{itemize}
		
		\vspace{0.1cm}
		\textbf{Example:}
		\begin{itemize}
			\item It's not illegal to abandon friends, but it's immoral
			\item Laws provide minimums; ethics provides ideals
		\end{itemize}
		
		\textcolor{myblue}{\textbf{Key Insight:}} Ethics \textcolor{myblue}{complements and informs} law - goes beyond minimum requirements!
	\end{frame}
	
	% ============================================================
	% SECTION 2: BENEFITS OF ETHICAL FRAMEWORKS - SLIDES 6-15
	% ============================================================
	
	\begin{frame}{Slide 6: Strengthening Organizational Culture}
		\fontsize{6.5}{7.5}\selectfont
		\textbf{How Ethics Strengthens Culture:}
		
		\vspace{0.1cm}
		\textbf{Key Cultural Benefits:}
		\begin{enumerate}
			\item \textbf{Teamwork:}
			\begin{itemize}
				\item Ethics programs nurture collaboration
				\item Shared values create alignment
				\item Better team dynamics
			\end{itemize}
			\item \textbf{Accountability:}
			\begin{itemize}
				\item Clear ethical standards
				\item Responsibility for actions
				\item Transparent decision-making
			\end{itemize}
			\item \textbf{Integrity:}
			\begin{itemize}
				\item Doing the right thing
				\item Ethical leadership
				\item Trustworthy practices
			\end{itemize}
		\end{enumerate}
		
		\vspace{0.1cm}
		\textbf{Results:}
		\begin{itemize}
			\item Boosts morale
			\item Increases productivity
			\item Creates positive work environment
			\item Reduces ethical lapses
		\end{itemize}
		
		\textcolor{myblue}{\textbf{Key Insight:}} Ethics frameworks strengthen \textcolor{myblue}{teamwork, accountability, and integrity}!
	\end{frame}
	
	\begin{frame}{7: Supporting Risk Management}
		\fontsize{6.5}{7.5}\selectfont
		\textbf{Risk Management Benefits of Ethics Frameworks:}
		
		\vspace{0.1cm}
		\textbf{Key Risk Reductions:}
		\begin{enumerate}
			\item \textbf{Safety Failures:}
			\begin{itemize}
				\item Ethical design considers safety
				\item Proactive risk identification
				\item Prevention of harmful outcomes
			\end{itemize}
			\item \textbf{Data Breaches:}
			\begin{itemize}
				\item Privacy as ethical priority
				\item Strong security measures
				\item Data minimization
			\end{itemize}
			\item \textbf{Regulatory Non-Compliance:}
			\begin{itemize}
				\item Ethics frameworks anticipate regulation
				\item Proactive compliance
				\item Avoid penalties
			\end{itemize}
		\end{enumerate}
		
		\vspace{0.1cm}
		\textbf{How It Works:}
		\begin{itemize}
			\item Ethical frameworks undergird policies and procedures
			\item Provide guidance for difficult decisions
			\item Create accountability structures
		\end{itemize}
		
		\textcolor{myblue}{\textbf{Key Insight:}} Ethics frameworks \textcolor{myblue}{reduce risks} of failures, breaches, and non-compliance!
	\end{frame}
	
	\begin{frame}{Slide 8: Guiding Responsible Innovation}
		\fontsize{6.5}{7.5}\selectfont
		\textbf{How Ethics Guides Innovation:}
		
		\vspace{0.1cm}
		\textbf{Responsible Innovation:}
		\begin{enumerate}
			\item \textbf{Consider Impacts:}
			\begin{itemize}
				\item Ethics ensures consideration of people and society
				\item Balances progress with protection
				\item Prevents harmful applications
			\end{itemize}
			\item \textbf{Direct Development:}
			\begin{itemize}
				\item Ethics provides direction for innovation
				\item Identifies acceptable boundaries
				\item Encourages beneficial applications
			\end{itemize}
			\item \textbf{Build Trust:}
			\begin{itemize}
				\item Ethical innovation builds public trust
				\item Encourages adoption
				\item Creates sustainable progress
			\end{itemize}
		\end{enumerate}
		
		\vspace{0.1cm}
		\textbf{Key Principle:}
		\begin{itemize}
			\item Innovation should benefit society
			\item Ethics ensures technology serves human values
			\item Prevents technology from serving only narrow interests
		\end{itemize}
		
		\textcolor{myblue}{\textbf{Key Insight:}} Ethics guides innovation in \textcolor{myblue}{responsible directions} considering societal impacts!
	\end{frame}
	
	\begin{frame}{Slide 9: Promoting Long-Term Thinking}
		\fontsize{6.5}{7.5}\selectfont
		\textbf{How Ethics Promotes Long-Term Perspectives:}
		
		\vspace{0.1cm}
		\textbf{Long-Term Benefits:}
		\begin{enumerate}
			\item \textbf{Sustainability:}
			\begin{itemize}
				\item Considering future impacts
				\item Environmental and social sustainability
				\item Intergenerational justice
			\end{itemize}
			\item \textbf{Stakeholder Value:}
			\begin{itemize}
				\item Balancing short and long-term interests
				\item Considering all affected parties
				\item Building lasting relationships
			\end{itemize}
			\item \textbf{Resilience:}
			\begin{itemize}
				\item Ethical companies weather crises better
				\item Trust provides buffer
				\item Long-term viability
			\end{itemize}
		\end{enumerate}
		
		\vspace{0.1cm}
		\textbf{Key Insight:}
		\begin{itemize}
			\item Managing with ethical considerations promotes long-term view
			\item Counteracts short-term profit pressures
			\item Builds enduring value
		\end{itemize}
		
		\textcolor{myblue}{\textbf{Key Insight:}} Ethics promotes \textcolor{myblue}{long-term thinking} and sustainability!
	\end{frame}
	
	\begin{frame}{Slide 10: Business Case for Ethics}
		\fontsize{6.5}{7.5}\selectfont
		\textbf{The Business Case for Practical Ethics:}
		
		\vspace{0.1cm}
		\textbf{Why Ethics Makes Business Sense:}
		\begin{enumerate}
			\item \textbf{Trust and Reputation:}
			\begin{itemize}
				\item Ethical companies are trusted
				\item Better brand image
				\item Customer loyalty
			\end{itemize}
			\item \textbf{Risk Reduction:}
			\begin{itemize}
				\item Avoid scandals
				\item Prevent legal issues
				\item Reduce regulatory exposure
			\end{itemize}
			\item \textbf{Talent Acquisition:}
			\begin{itemize}
				\item Attract values-aligned employees
				\item Retain skilled workers
				\item Build committed teams
			\end{itemize}
			\item \textbf{Culture:}
			\begin{itemize}
				\item Stronger workplace culture
				\item Higher productivity
				\item Better collaboration
			\end{itemize}
		\end{enumerate}
		
		\vspace{0.1cm}
		\textbf{Key Point:}
		Proponents argue that practical ethics is not just morally right but also financially prudent.
		
		\textcolor{myblue}{\textbf{Key Insight:}} Ethics makes \textcolor{myblue}{good business sense} with measurable benefits!
	\end{frame}
	
	\begin{frame}{Slide 11: Ethics in the Workplace}
		\fontsize{6.5}{7.5}\selectfont
		\textbf{How Ethics Manifests in Organizations:}
		
		\vspace{0.1cm}
		\textbf{Practical Applications:}
		\begin{enumerate}
			\item \textbf{Decision-Making:}
			\begin{itemize}
				\item Ethics provides frameworks for difficult choices
				\item Guides trade-off decisions
				\item Ensures consistency
			\end{itemize}
			\item \textbf{Policy Development:}
			\begin{itemize}
				\item Ethics undergirds policies
				\item Creates accountability structures
				\item Establishes standards
			\end{itemize}
			\item \textbf{Training:}
			\begin{itemize}
				\item Ethics training for employees
				\item Awareness of ethical issues
				\item Skills for ethical decision-making
			\end{itemize}
			\item \textbf{Oversight:}
			\begin{itemize}
				\item Ethics committees
				\item Review processes
				\item Monitoring and enforcement
			\end{itemize}
		\end{enumerate}
		
		\textcolor{myblue}{\textbf{Key Insight:}} Ethics must be \textcolor{myblue}{embedded in organizational practices}!
	\end{frame}
	
	\begin{frame}{Slide 12: Regulatory Recognition}
		\fontsize{6.5}{7.5}\selectfont
		\textbf{Regulatory Recognition of Ethics:}
		
		\vspace{0.1cm}
		\textbf{Why Regulators Care About Ethics:}
		\begin{enumerate}
			\item \textbf{Protecting Citizens:}
			\begin{itemize}
				\item Ethical standards protect individuals
				\item Prevent harm and discrimination
				\item Ensure accountability
			\end{itemize}
			\item \textbf{Market Confidence:}
			\begin{itemize}
				\item Ethical practices build trust
				\item Stable markets require trust
				\item Investor confidence
			\end{itemize}
			\item \textbf{International Standards:}
			\begin{itemize}
				\item Ethics frameworks guide regulation
				\item Global alignment
				\item Harmonized standards
			\end{itemize}
		\end{enumerate}
		
		\vspace{0.1cm}
		\textbf{Key Trend:}
		Regulations increasingly require ethical considerations (e.g., GDPR, EU AI Act).
		
		\textcolor{myblue}{\textbf{Key Insight:}} Ethics frameworks underlay \textcolor{myblue}{regulatory and legislative approaches}!
	\end{frame}
	
	\begin{frame}{Slide 13: Ethics vs Compliance}
		\fontsize{6.5}{7.5}\selectfont
		\textbf{Understanding Ethics vs Compliance:}
		
		\vspace{0.1cm}
		\textbf{Compliance:}
		\begin{itemize}
			\item Following laws and regulations
			\item Minimum acceptable behavior
			\item External requirement
			\item Reactive approach
			\item "What must we do?"
		\end{itemize}
		
		\vspace{0.1cm}
		\textbf{Ethics:}
		\begin{itemize}
			\item Going beyond legal requirements
			\item Aspirational standards
			\item Internal commitment
			\item Proactive approach
			\item "What should we do?"
		\end{itemize}
		
		\vspace{0.1cm}
		\textbf{Key Insight:}
		\begin{itemize}
			\item Compliance is necessary but not sufficient
			\item Ethics provides the moral compass
			\item Ethics prevents issues that compliance doesn't address
		\end{itemize}
		
		\textcolor{myblue}{\textbf{Key Insight:}} Ethics goes \textcolor{myblue}{beyond compliance} - establishes aspirational standards!
	\end{frame}
	
	\begin{frame}{Slide 14: Stakeholder Expectations}
		\fontsize{6.5}{7.5}\selectfont
		\textbf{Stakeholder Expectations for Ethical AI:}
		
		\vspace{0.1cm}
		\textbf{Key Stakeholder Groups:}
		\begin{enumerate}
			\item \textbf{Customers:}
			\begin{itemize}
				\item Expect fair treatment
				\item Privacy protection
				\item Transparent decision-making
			\end{itemize}
			\item \textbf{Employees:}
			\begin{itemize}
				\item Ethical workplace
				\item Fair algorithms
				\item Meaningful work
			\end{itemize}
			\item \textbf{Investors:}
			\begin{itemize}
				\item Sustainable practices
				\item Risk management
				\item Long-term value
			\end{itemize}
			\item \textbf{Regulators:}
			\begin{itemize}
				\item Compliance
				\item Accountability
				\item Transparency
			\end{itemize}
			\item \textbf{Society:}
			\begin{itemize}
				\item Social good
				\item Fairness
				\item Trust
			\end{itemize}
		\end{enumerate}
		
		\textcolor{myblue}{\textbf{Key Insight:}} All stakeholders expect \textcolor{myblue}{ethical AI practices}!
	\end{frame}
	
	\begin{frame}{Slide 15: Ethics Benefits Summary}
		\fontsize{6.5}{7.5}\selectfont
		\textbf{Comprehensive Benefits of Practical Ethics:}
		
		\vspace{0.1cm}
		\textbf{Summary of Key Benefits:}
		\begin{enumerate}
			\item \textbf{Trust and Reputation:}
			\begin{itemize}
				\item Enhanced brand image
				\item Customer loyalty
				\item Stakeholder confidence
			\end{itemize}
			\item \textbf{Risk Management:}
			\begin{itemize}
				\item Avoid scandals
				\item Reduce legal exposure
				\item Prevent regulatory issues
			\end{itemize}
			\item \textbf{Talent Management:}
			\begin{itemize}
				\item Attract top talent
				\item Retain employees
				\item Build engaged workforce
			\end{itemize}
			\item \textbf{Culture:}
			\begin{itemize}
				\item Stronger organizational culture
				\item Better teamwork
				\item Higher productivity
			\end{itemize}
			\item \textbf{Innovation:}
			\begin{itemize}
				\item Responsible innovation
				\item Sustainable development
				\item Long-term thinking
			\end{itemize}
		\end{enumerate}
		
		\textcolor{myblue}{\textbf{Key Insight:}} Ethics provides \textcolor{myblue}{multiple overlapping benefits} to organizations!
	\end{frame}
	
	% ============================================================
	% SECTION 3: ETHICAL FRAMEWORKS - SLIDES 16-35
	% ============================================================
	
	\begin{frame}{Slide 16: Overview of Ethical Frameworks}
		\fontsize{6.5}{7.5}\selectfont
		\textbf{Three Major Ethical Frameworks:}
		
		\vspace{0.1cm}
		\textbf{1. Consequentialism:}
		\begin{itemize}
			\item Judges morality based on consequences
			\item Most common: Utilitarianism
			\item Focus: Maximize overall utility
			\item Forward-looking, outcome-focused
		\end{itemize}
		
		\vspace{0.1cm}
		\textbf{2. Deontology:}
		\begin{itemize}
			\item Judges morality based on duties and rules
			\item Key thinker: Immanuel Kant
			\item Focus: Adherence to ethical principles
			\item Duty-based, rule-focused
		\end{itemize}
		
		\vspace{0.1cm}
		\textbf{3. Virtue Ethics:}
		\begin{itemize}
			\item Emphasizes virtuous character traits
			\item Key thinker: Aristotle
			\item Focus: Living a good life
			\item Character-based, holistic
		\end{itemize}
		
		\vspace{0.1cm}
		\textbf{Key Insight:}
		These frameworks complement one another in practical ethics.
		
		\textcolor{myblue}{\textbf{Key Insight:}} Three frameworks offer \textcolor{myblue}{different but complementary perspectives}!
	\end{frame}
	
	\begin{frame}{Slide 17: Consequentialism Introduction}
		\fontsize{6.5}{7.5}\selectfont
		\textbf{Consequentialism Explained:}
		
		\vspace{0.1cm}
		\textbf{Core Idea:}
		\begin{itemize}
			\item Actions are judged by their consequences
			\item The morality of an action depends on its outcomes
			\item "The ends justify the means" (in some interpretations)
		\end{itemize}
		
		\vspace{0.1cm}
		\textbf{Key Principles:}
		\begin{enumerate}
			\item \textbf{Outcome Focus:}
			\begin{itemize}
				\item What results from the action?
				\item Does it produce good or bad outcomes?
				\item Consequences determine moral value
			\end{itemize}
			\item \textbf{Forward-Looking:}
			\begin{itemize}
				\item Focus on future results
				\item Past actions judged by their effects
				\item Decisions based on expected outcomes
			\end{itemize}
			\item \textbf{Circumstantially Relative:}
			\begin{itemize}
				\item What's right depends on circumstances
				\item No absolute rules
				\item Context matters
			\end{itemize}
		\end{enumerate}
		
		\textcolor{myblue}{\textbf{Key Insight:}} Consequentialism = \textcolor{myblue}{judging actions by their consequences}!
	\end{frame}
	
	\begin{frame}{Slide 18: Utilitarianism}
		\fontsize{6.5}{7.5}\selectfont
		\textbf{Utilitarianism: The Most Common Form of Consequentialism}
		
		\vspace{0.1cm}
		\textbf{Core Idea:}
		\begin{itemize}
			\item Maximize overall utility
			\item Utility often defined as happiness, pleasure, or satisfaction
			\item The morally correct action produces the greatest net utility
		\end{itemize}
		
		\vspace{0.1cm}
		\textbf{Key Thinkers:}
		\begin{enumerate}
			\item \textbf{Jeremy Bentham:}
			\begin{itemize}
				\item Founder of utilitarianism
				\item Hedonic calculus
				\item Focus on pleasure and pain
			\end{itemize}
			\item \textbf{John Stuart Mill:}
			\begin{itemize}
				\item Refined utilitarianism
				\item Quality of pleasures
				\item Higher and lower pleasures
			\end{itemize}
		\end{enumerate}
		
		\vspace{0.1cm}
		\textbf{Application to AI:}
		\begin{itemize}
			\item An AI system is morally good if it increases overall pleasure/happiness
			\item Actions are not intrinsically moral
			\item Derive moral value solely from results
		\end{itemize}
		
		\textcolor{myblue}{\textbf{Key Insight:}} Utilitarianism = \textcolor{myblue}{maximize overall utility/happiness}!
	\end{frame}
	
	\begin{frame}{Slide 19: Consequentialism Advantages}
		\fontsize{6.5}{7.5}\selectfont
		\textbf{Advantages of Consequentialism:}
		
		\vspace{0.1cm}
		\textbf{Key Advantages:}
		\begin{enumerate}
			\item \textbf{Single Quantifiable Metric:}
			\begin{itemize}
				\item Provides one measure of moral value
				\item Utility can be calculated
				\item Clear decision-making framework
			\end{itemize}
			\item \textbf{Impartiality:}
			\begin{itemize}
				\item Everyone's utility counts equally
				\item No special treatment
				\item Objective calculation
			\end{itemize}
			\item \textbf{Scientific Amenability:}
			\begin{itemize}
				\item Utility can be measured
				\item Evidence-based decisions
				\item Quantifiable outcomes
			\end{itemize}
			\item \textbf{Forward-Looking:}
			\begin{itemize}
				\item Focuses on future outcomes
				\item Practical decision-making
				\item Real-world consequences
			\end{itemize}
		\end{enumerate}
		
		\vspace{0.1cm}
		\textbf{Note:}
		Quantification of value can be challenging despite these advantages.
		
		\textcolor{myblue}{\textbf{Key Insight:}} Consequentialism provides \textcolor{myblue}{a single quantifiable metric} for moral decisions!
	\end{frame}
	
	\begin{frame}{Slide 20: Consequentialism Challenges}
		\fontsize{6.5}{7.5}\selectfont
		\textbf{Criticisms and Challenges of Consequentialism:}
		
		\vspace{0.1cm}
		\textbf{Key Challenges:}
		\begin{enumerate}
			\item \textbf{Subjectivity of Utility:}
			\begin{itemize}
				\item Choice of metric is subjective
				\item Different definitions of "good"
				\item Measurement difficulties
			\end{itemize}
			\item \textbf{Unpredictable Consequences:}
			\begin{itemize}
				\item Can't predict all outcomes
				\item Where does the calculation stop?
				\item Immediate vs. long-term consequences
			\end{itemize}
			\item \textbf{Rights Violations:}
			\begin{itemize}
				\item Can justify violating individual rights
				\item "The ends justify the means"
				\item Classic example: Killing one to save five
			\end{itemize}
			\item \textbf{Doctor Example:}
			\begin{itemize}
				\item Killing a patient to use organs saves five
				\item Maximizes utility but seems immoral
				\item Challenge for consequentialism
			\end{itemize}
		\end{enumerate}
		
		\textcolor{myblue}{\textbf{Key Insight:}} Consequentialism can justify \textcolor{myblue}{violating rights for the greater good}!
	\end{frame}
	
	\begin{frame}{Slide 21: Rule Consequentialism}
		\fontsize{6.5}{7.5}\selectfont
		\textbf{Rule Consequentialism: A Modification}
		
		\vspace{0.1cm}
		\textbf{Core Idea:}
		\begin{itemize}
			\item Judges acts by whether they adhere to utility-maximizing rules
			\item Not just by case-specific outcomes
			\item Addresses issues with act utilitarianism
		\end{itemize}
		
		\vspace{0.1cm}
		\textbf{How It Works:}
		\begin{enumerate}
			\item \textbf{Follow General Rules:}
			\begin{itemize}
				\item Rules that tend to maximize utility
				\item Rules against murder, lying, etc.
				\item Upheld even when breaking them may increase utility
			\end{itemize}
			\item \textbf{Example - Doctor Case:}
			\begin{itemize}
				\item Doctor should not kill patient
				\item Would undermine trust in medical system
				\item Long-term utility is maximized by following rules
			\end{itemize}
			\item \textbf{AI Application:}
			\begin{itemize}
				\item Should not violate rights even if efficient
				\item Long-term trust is valuable
				\item Rule-based protection
			\end{itemize}
		\end{enumerate}
		
		\textcolor{myblue}{\textbf{Key Insight:}} Rule consequentialism = \textcolor{myblue}{follow utility-maximizing rules}!
	\end{frame}
	
	\begin{frame}{Slide 22: Deontology Introduction}
		\fontsize{6.5}{7.5}\selectfont
		\textbf{Deontology Explained:}
		
		\vspace{0.1cm}
		\textbf{Core Idea:}
		\begin{itemize}
			\item Judges morality based on adherence to duties and rules
			\item Not based on consequences
			\item Focus on moral obligations regardless of outcomes
		\end{itemize}
		
		\vspace{0.1cm}
		\textbf{Key Principles:}
		\begin{enumerate}
			\item \textbf{Duty-Based:}
			\begin{itemize}
				\item Moral obligations are paramount
				\item We must do our duty
				\item Consequences don't override duties
			\end{itemize}
			\item \textbf{Rule-Focused:}
			\begin{itemize}
				\item Follow moral rules
				\item Universal principles
				\item Consistency across situations
			\end{itemize}
			\item \textbf{Respect for Persons:}
			\begin{itemize}
				\item People are ends in themselves
				\item Never treat as mere means
				\item Intrinsic value of individuals
			\end{itemize}
		\end{enumerate}
		
		\textcolor{myblue}{\textbf{Key Insight:}} Deontology = \textcolor{myblue}{judging actions by duties and rules}!
	\end{frame}
	
	\begin{frame}{Slide 23: Immanuel Kant and Deontology}
		\fontsize{6.5}{7.5}\selectfont
		\textbf{Kant's Contribution to Deontology:}
		
		\vspace{0.1cm}
		\textbf{Key Thinker:}
		\begin{itemize}
			\item Immanuel Kant is the most important proponent of deontology
			\item Emphasized moral principles, rationality, and duty
			\item Famous for not placing moral weight on consequences
			\item Believed lying is always wrong regardless of consequences
		\end{itemize}
		
		\vspace{0.1cm}
		\textbf{Core Concepts:}
		\begin{enumerate}
			\item \textbf{Categorical Imperative:}
			\begin{itemize}
				\item Act only according to maxims that could be universal law
				\item "It's wrong to lie" - applicable universally
				\item No contradictions in universal application
			\end{itemize}
			\item \textbf{Treating People as Ends:}
			\begin{itemize}
				\item Never treat people as mere means to an end
				\item People have intrinsic value
				\item Respect their autonomy and dignity
			\end{itemize}
			\item \textbf{Reason-Based Morality:}
			\begin{itemize}
				\item Morality grounded in reason
				\item Not in emotions or consequences
				\item Rational beings can determine moral duties
			\end{itemize}
		\end{enumerate}
		
		\textcolor{myblue}{\textbf{Key Insight:}} Kant = \textcolor{myblue}{duty-based ethics focused on universal principles}!
	\end{frame}
	
	\begin{frame}{Slide 24: The Categorical Imperative}
		\fontsize{6.5}{7.5}\selectfont
		\textbf{Understanding the Categorical Imperative:}
		
		\vspace{0.1cm}
		\textbf{First Formulation:}
		\begin{itemize}
			\item "Act only according to that maxim whereby you can at the same time will that it should become a universal law"
			\item Can the rule be applied universally without contradiction?
			\item Example: Could "It's okay to lie" be a universal law?
		\end{itemize}
		
		\vspace{0.1cm}
		\textbf{Second Formulation:}
		\begin{itemize}
			\item "Act in such a way that you treat humanity, whether in your own person or in the person of any other, always at the same time as an end and never merely as a means"
			\item Respect the dignity of all persons
			\item Don't use people as tools
		\end{itemize}
		
		\vspace{0.1cm}
		\textbf{Application to AI:}
		\begin{itemize}
			\item AI systems should respect human dignity
			\item People should not be treated as data points
			\item Transparency and consent are essential
		\end{itemize}
		
		\textcolor{myblue}{\textbf{Key Insight:}} Categorical imperative = \textcolor{myblue}{act on rules that could be universal law}!
	\end{frame}
	
	\begin{frame}{Slide 25: Modern Deontology}
		\fontsize{6.5}{7.5}\selectfont
		\textbf{Contemporary Deontological Ethics:}
		
		\vspace{0.1cm}
		\textbf{Evolution of Deontology:}
		\begin{itemize}
			\item Most contemporary deontologists care about consequences
			\item However, consequences are not the only moral consideration
			\item There are red lines that should not be crossed
		\end{itemize}
		
		\vspace{0.1cm}
		\textbf{Key Features:}
		\begin{enumerate}
			\item \textbf{Rights:}
			\begin{itemize}
				\item Individuals have fundamental rights
				\item Rights cannot be violated even for good outcomes
				\item Examples: Privacy, autonomy, non-discrimination
			\end{itemize}
			\item \textbf{Rules:}
			\begin{itemize}
				\item Established ethical rules matter
				\item Professional codes of conduct
				\item Organizational policies
			\end{itemize}
			\item \textbf{Principles:}
			\begin{itemize}
				\item Ethical principles provide guidance
				\item Beyond consequences
				\item Deontological in nature
			\end{itemize}
		\end{enumerate}
		
		\textcolor{myblue}{\textbf{Key Insight:}} Modern deontology \textcolor{myblue}{considers consequences but has red lines}!
	\end{frame}
	
	\begin{frame}{Slide 26: Deontology Challenges}
		\fontsize{6.5}{7.5}\selectfont
		\textbf{Criticisms and Challenges of Deontology:}
		
		\vspace{0.1cm}
		\textbf{Key Challenges:}
		\begin{enumerate}
			\item \textbf{Rigidity:}
			\begin{itemize}
				\item Can lead to rigid rule-following
				\item Even when consequences are disastrous
				\item Example: Kant said lying is always wrong
			\end{itemize}
			\item \textbf{Conflicting Duties:}
			\begin{itemize}
				\item What if duties conflict?
				\item Example: Duty to tell truth vs. duty to protect
				\item No clear resolution framework
			\end{itemize}
			\item \textbf{Application Difficulty:}
			\begin{itemize}
				\item Hard to apply in complex situations
				\item What rules apply?
				\item How to resolve conflicts?
			\end{itemize}
		\end{enumerate}
		
		\vspace{0.1cm}
		\textbf{AI Application:}
		\begin{itemize}
			\item Can lead to overly rigid AI systems
			\item May not adapt to context
			\item Difficulty encoding all rules
		\end{itemize}
		
		\textcolor{myblue}{\textbf{Key Insight:}} Deontology can be \textcolor{myblue}{too rigid} even when consequences are disastrous!
	\end{frame}
	
	\begin{frame}{Slide 27: Virtue Ethics Introduction}
		\fontsize{6.5}{7.5}\selectfont
		\textbf{Virtue Ethics Explained:}
		
		\vspace{0.1cm}
		\textbf{Core Idea:}
		\begin{itemize}
			\item Emphasizes virtuous character traits
			\item Focus on living a good life
			\item Rather than rules or consequences
		\end{itemize}
		
		\vspace{0.1cm}
		\textbf{Key Principles:}
		\begin{enumerate}
			\item \textbf{Character Focus:}
			\begin{itemize}
				\item What kind of person should I be?
				\item Cultivating virtues
				\item Becoming a good person
			\end{itemize}
			\item \textbf{Virtues:}
			\begin{itemize}
				\item Wisdom, courage, humanity
				\item Justice, temperance, generosity
				\item Character traits to cultivate
			\end{itemize}
			\item \textbf{Practical Wisdom:}
			\begin{itemize}
				\item Exercising judgment in situations
				\item Moderating emotions and appetites
				\item Appropriate behavior in context
			\end{itemize}
		\end{enumerate}
		
		\textcolor{myblue}{\textbf{Key Insight:}} Virtue ethics = \textcolor{myblue}{focus on character and living a good life}!
	\end{frame}
	
	\begin{frame}{Slide 28: Aristotle and Virtue Ethics}
		\fontsize{6.5}{7.5}\selectfont
		\textbf{Aristotle's Contribution:}
		
		\vspace{0.1cm}
		\textbf{Historical Roots:}
		\begin{itemize}
			\item Virtue ethics has roots in ancient Greek philosophers
			\item Aristotle taught that happiness comes from a life guided by virtues
			\item Virtues are nurtured through practice and habit
		\end{itemize}
		
		\vspace{0.1cm}
		\textbf{Key Concepts:}
		\begin{enumerate}
			\item \textbf{Eudaimonia:}
			\begin{itemize}
				\item Flourishing or living well
				\item Ultimate goal of human life
				\item Achieved through virtue
			\end{itemize}
			\item \textbf{Virtues as Habits:}
			\begin{itemize}
				\item Virtues are developed through practice
				\item Repeated actions create character
				\item Modeling virtuous exemplars
			\end{itemize}
			\item \textbf{Golden Mean:}
			\begin{itemize}
				\item Virtue lies between extremes
				\item Courage between cowardice and recklessness
				\item Balance in character
			\end{itemize}
		\end{enumerate}
		
		\textcolor{myblue}{\textbf{Key Insight:}} Aristotle = \textcolor{myblue}{virtues through practice for a good life}!
	\end{frame}
	
	\begin{frame}{Slide 29: Key Virtues}
		\fontsize{6.5}{7.5}\selectfont
		\textbf{Examples of Virtues:}
		
		\vspace{0.1cm}
		\textbf{Core Virtues:}
		\begin{enumerate}
			\item \textbf{Wisdom:}
			\begin{itemize}
				\item Practical judgment
				\item Understanding what matters
				\item Sound decision-making
			\end{itemize}
			\item \textbf{Courage:}
			\begin{itemize}
				\item Facing challenges appropriately
				\item Moral courage
				\item Standing up for what's right
			\end{itemize}
			\item \textbf{Humanity:}
			\begin{itemize}
				\item Compassion and care
				\item Empathy for others
				\item Benevolence
			\end{itemize}
			\item \textbf{Justice:}
			\begin{itemize}
				\item Fairness and equity
				\item Giving proper due
				\item Social responsibility
			\end{itemize}
			\item \textbf{Temperance:}
			\begin{itemize}
				\item Self-control
				\item Moderation
				\item Balance
			\end{itemize}
			\item \textbf{Generosity:}
			\begin{itemize}
				\item Giving freely
				\item Sharing resources
				\item Beneficence
			\end{itemize}
		\end{enumerate}
		
		\textcolor{myblue}{\textbf{Key Insight:}} Virtues include \textcolor{myblue}{wisdom, courage, justice, temperance, and generosity}!
	\end{frame}
	
	\begin{frame}{Slide 30: Virtue Ethics in Practice}
		\fontsize{6.5}{7.5}\selectfont
		\textbf{Applying Virtue Ethics:}
		
		\vspace{0.1cm}
		\textbf{Practical Approach:}
		\begin{enumerate}
			\item \textbf{Ask the Virtue Question:}
			\begin{itemize}
				\item "What would a virtuous person do?"
				\item Consider role models and exemplars
				\item What would Jesus, Mohammed, Gandhi do?
			\end{itemize}
			\item \textbf{Cultivate Virtues:}
			\begin{itemize}
				\item Through practice and habit
				\item Learning from experience
				\item Modeling virtuous behavior
			\end{itemize}
			\item \textbf{Consider Character:}
			\begin{itemize}
				\item Not just discrete acts
				\item Life-long character development
				\item Holistic approach to morality
			\end{itemize}
		\end{enumerate}
		
		\vspace{0.1cm}
		\textbf{Application to AI:}
		\begin{itemize}
			\item Build AI that learns from experience like children
			\item Develop virtuous organizational culture
			\item Consider what responsible companies would do
		\end{itemize}
		
		\textcolor{myblue}{\textbf{Key Insight:}} Virtue ethics asks \textcolor{myblue}{"What would a virtuous person do?"}!
	\end{frame}
	
	\begin{frame}{Slide 31: Virtue Ethics Challenges}
		\fontsize{6.5}{7.5}\selectfont
		\textbf{Criticisms of Virtue Ethics:}
		
		\vspace{0.1cm}
		\textbf{Key Challenges:}
		\begin{enumerate}
			\item \textbf{Lack of Clear Guidance:}
			\begin{itemize}
				\item Less concrete than rules-based approaches
				\item Hard to know what virtue requires
				\item Situational judgment needed
			\end{itemize}
			\item \textbf{Conflicting Virtues:}
			\begin{itemize}
				\item Different virtues can conflict
				\item How to weigh courage vs. prudence?
				\item No clear resolution method
			\end{itemize}
			\item \textbf{Cultural Relativism:}
			\begin{itemize}
				\item Virtues vary across cultures
				\item What's virtuous in one context may not be in another
				\item Need for universal standards
			\end{itemize}
		\end{enumerate}
		
		\vspace{0.1cm}
		\textbf{Virtue Ethicist Response:}
		\begin{itemize}
			\item Practical wisdom helps navigate hard cases
			\item Other theories also have conflicts
			\item Consequences not always comparable
		\end{itemize}
		
		\textcolor{myblue}{\textbf{Key Insight:}} Virtue ethics criticized for \textcolor{myblue}{lack of clear guidance}!
	\end{frame}
	
	\begin{frame}{Slide 32: Modern Virtue Ethics}
		\fontsize{6.5}{7.5}\selectfont
		\textbf{Contemporary Applications of Virtue Ethics:}
		
		\vspace{0.1cm}
		\textbf{Extension to Organizations:}
		\begin{itemize}
			\item Virtue ethics expands beyond individual character
			\item Organizational virtues include:
			\begin{itemize}
				\item Justice
				\item Accountability
				\item Environmental stewardship
				\item Responsible innovation
			\end{itemize}
		\end{itemize}
		
		\vspace{0.1cm}
		\textbf{Renewed Interest:}
		\begin{itemize}
			\item Across moral psychology
			\item Business ethics
			\item AI ethics and development
		\end{itemize}
		
		\vspace{0.1cm}
		\textbf{Integration with Common Morality:}
		\begin{itemize}
			\item People learn morality through habituation
			\item Parents teach children through repeated examples
			\item Virtue ethics aligns with how morality develops
		\end{itemize}
		
		\textcolor{myblue}{\textbf{Key Insight:}} Organizational virtues include \textcolor{myblue}{justice, accountability, and responsible innovation}!
	\end{frame}
	
	\begin{frame}{Slide 33: Virtue Ethics and AI}
		\fontsize{6.5}{7.5}\selectfont
		\textbf{Applying Virtue Ethics to AI:}
		
		\vspace{0.1cm}
		\textbf{Virtue-Based AI Development:}
		\begin{enumerate}
			\item \textbf{Learning from Experience:}
			\begin{itemize}
				\item Some scholars suggest AI should learn like children
				\item Through experience and habit
				\item Not just coded rules
			\end{itemize}
			\item \textbf{Organizational Virtues:}
			\begin{itemize}
				\item Responsible development
				\item Ethical corporate culture
				\item Virtuous design teams
			\end{itemize}
			\item \textbf{Character of AI:}
			\begin{itemize}
				\item What would an ethical AI do?
				\item Embedding virtues in AI systems
				\item Training on ethical examples
			\end{itemize}
		\end{enumerate}
		
		\vspace{0.1cm}
		\textbf{Key Insight:}
		\begin{itemize}
			\item Morality is complex - may be impossible to code top-down
			\item Learning approach may be more effective
			\item Ethics as character development
		\end{itemize}
		
		\textcolor{myblue}{\textbf{Key Insight:}} Virtue ethics AI would \textcolor{myblue}{learn from experience like children}!
	\end{frame}
	
	\begin{frame}{Slide 34: Framework Complementarity}
		\fontsize{6.5}{7.5}\selectfont
		\textbf{How Frameworks Complement Each Other:}
		
		\vspace{0.1cm}
		\textbf{Combined Approach:}
		\begin{itemize}
			\item Consequentialism, deontology, and virtue ethics are not mutually exclusive
			\item In practical ethics, they complement one another
			\item The best moral decision accords with all three
		\end{itemize}
		
		\vspace{0.1cm}
		\textbf{Integrated Decision-Making:}
		\begin{enumerate}
			\item \textbf{Consequentialist Analysis:}
			\begin{itemize}
				\item What are the likely outcomes?
				\item Does it maximize good consequences?
				\item Consider all affected parties
			\end{itemize}
			\item \textbf{Deontological Analysis:}
			\begin{itemize}
				\item Does it respect rights and duties?
				\item Are there red lines being crossed?
				\item Is it in accordance with principles?
			\end{itemize}
			\item \textbf{Virtue Ethics Analysis:}
			\begin{itemize}
				\item What would a virtuous person do?
				\item Does it build character?
				\item Is it consistent with values?
			\end{itemize}
		\end{enumerate}
		
		\textcolor{myblue}{\textbf{Key Insight:}} Frameworks are \textcolor{myblue}{complementary, not mutually exclusive}!
	\end{frame}
	
	\begin{frame}{Slide 35: Ethical Frameworks Summary}
		\fontsize{6.5}{7.5}\selectfont
		\textbf{Summary of Ethical Frameworks:}
		
		\vspace{0.1cm}
		\textbf{Consequentialism:}
		\begin{itemize}
			\item Focus: Outcomes and consequences
			\item Key Question: "What maximizes good outcomes?"
			\item Strength: Single quantifiable metric
			\item Weakness: Can justify rights violations
		\end{itemize}
		
		\vspace{0.1cm}
		\textbf{Deontology:}
		\begin{itemize}
			\item Focus: Duties and rules
			\item Key Question: "What are our moral obligations?"
			\item Strength: Respects rights and dignity
			\item Weakness: Can be rigid
		\end{itemize}
		
		\vspace{0.1cm}
		\textbf{Virtue Ethics:}
		\begin{itemize}
			\item Focus: Character and virtues
			\item Key Question: "What kind of person should I be?"
			\item Strength: Holistic and aspirational
			\item Weakness: Lack of clear guidance
		\end{itemize}
		
		\vspace{0.1cm}
		\textbf{Key Insight:}
		\begin{itemize}
			\item Each framework offers valuable insights
			\item Best decisions often integrate all three
			\item Context determines emphasis
		\end{itemize}
		
		\textcolor{myblue}{\textbf{Key Insight:}} Three frameworks provide \textcolor{myblue}{complementary perspectives} for ethical AI!
	\end{frame}
	
	% ============================================================
	% SECTION 4: MEDICAL ETHICS AND AI - SLIDES 36-40
	% ============================================================
	
	\begin{frame}{Slide 36: Medical Ethics History}
		\fontsize{6.5}{7.5}\selectfont
		\textbf{Medical Ethics: A Historical Overview}
		
		\vspace{0.1cm}
		\textbf{Ancient Foundations:}
		\begin{itemize}
			\item Hippocratic Oath (5th-3rd century BC)
			\item "Do no harm" to patients
			\item First known code of medical ethics
		\end{itemize}
		
		\vspace{0.1cm}
		\textbf{Key Documents:}
		\begin{enumerate}
			\item \textbf{Nuremberg Code (1947):}
			\begin{itemize}
				\item Response to Nazi medical experiments
				\item Informed consent requirement
				\item Research ethics foundation
			\end{itemize}
			\item \textbf{Declaration of Geneva (1948):}
			\begin{itemize}
				\item Modern Hippocratic Oath
				\item International medical ethics
			\end{itemize}
			\item \textbf{Declaration of Helsinki (1964):}
			\begin{itemize}
				\item Research ethics guidelines
				\item Updated regularly
				\item Global standards
			\end{itemize}
			\item \textbf{Belmont Report (1978):}
			\begin{itemize}
				\item Respect for persons, beneficence, justice
				\item Research ethics framework
				\item US foundational document
			\end{itemize}
		\end{enumerate}
		
		\textcolor{myblue}{\textbf{Key Insight:}} Medical ethics has \textcolor{myblue}{long history and established frameworks}!
	\end{frame}
	
	\begin{frame}{Slide 37: Medical Ethics Evolution}
		\fontsize{6.5}{7.5}\selectfont
		\textbf{Medical Ethics Development in the 1970s:}
		
		\vspace{0.1cm}
		\textbf{Factors Driving Evolution:}
		\begin{enumerate}
			\item \textbf{Concentration of Medical Care:}
			\begin{itemize}
				\item Hospitals and depersonalized settings
				\item Loss of personal relationships
				\item Need for formal ethics
			\end{itemize}
			\item \textbf{Rising Costs and Government Role:}
			\begin{itemize}
				\item Increased government involvement
				\item Healthcare funding changes
				\item Resource allocation decisions
			\end{itemize}
			\item \textbf{Patient Rights Movements:}
			\begin{itemize}
				\item Civil rights outgrowth
				\item Demands for autonomy
				\item Informed consent
			\end{itemize}
			\item \textbf{Medical Scandals:}
			\begin{itemize}
				\item Tuskegee Syphilis Experiment
				\item Public outrage
				\item Call for reform
			\end{itemize}
			\item \textbf{Rapid Technological Advances:}
			\begin{itemize}
				\item Mechanical ventilators
				\item Organ transplantation
				\item New ethical dilemmas
			\end{itemize}
		\end{enumerate}
		
		\textcolor{myblue}{\textbf{Key Insight:}} Multiple factors drove \textcolor{myblue}{development of medical ethics}!
	\end{frame}
	
	\begin{frame}{Slide 38: Technological Advances and Ethics}
		\fontsize{6.5}{7.5}\selectfont
		\textbf{Technology Created New Ethical Challenges:}
		
		\vspace{0.1cm}
		\textbf{Example - Mechanical Ventilator:}
		\begin{itemize}
			\item Prompted reconsideration of death
			\item Heart-beating bodies with non-functioning brains
			\item Organ procurement dilemmas
			\item Is the patient alive or dead?
		\end{itemize}
		
		\vspace{0.1cm}
		\textbf{Key Lesson:}
		\begin{itemize}
			\item Technological advances create ethical challenges
			\item Moral questions, not medical ones
			\item Need for ethical frameworks
		\end{itemize}
		
		\vspace{0.1cm}
		\textbf{Parallels to AI:}
		\begin{itemize}
			\item AI creates new ethical dilemmas
			\item Not purely technical questions
			\item Need for ethical guidance
		\end{itemize}
		
		\textcolor{myblue}{\textbf{Key Insight:}} Technology creates \textcolor{myblue}{new ethical dilemmas} that need frameworks!
	\end{frame}
	
	\begin{frame}{Slide 39: Lessons for AI Ethics}
		\fontsize{6.5}{7.5}\selectfont
		\textbf{What AI Ethics Can Learn from Medical Ethics:}
		
		\vspace{0.1cm}
		\textbf{Key Lessons:}
		\begin{enumerate}
			\item \textbf{Structured Frameworks:}
			\begin{itemize}
				\item Medical ethics has established principles
				\item AI needs similar foundations
				\item Clear ethical guidance
			\end{itemize}
			\item \textbf{Professional Ethics:}
			\begin{itemize}
				\item Doctors have ethical training
				\item AI professionals need ethics education
				\item Professional responsibility
			\end{itemize}
			\item \textbf{Ethics Committees:}
			\begin{itemize}
				\item Medical ethics committees
				\item AI needs similar oversight
				\item Multi-disciplinary review
			\end{itemize}
			\item \textbf{Codes of Ethics:}
			\begin{itemize}
				\item Medical codes are well-established
				\item AI codes are evolving
				\item Need for standards
			\end{itemize}
		\end{enumerate}
		
		\textcolor{myblue}{\textbf{Key Insight:}} Medical ethics provides \textcolor{myblue}{model for AI ethics development}!
	\end{frame}
	
	\begin{frame}{Slide 40: Medical Ethics and AI Parallels}
		\fontsize{6.5}{7.5}\selectfont
		\textbf{Parallels Between Medical and AI Ethics:}
		
		\vspace{0.1cm}
		\textbf{Similar Challenges:}
		\begin{enumerate}
			\item \textbf{High-Stakes Decisions:}
			\begin{itemize}
				\item Medicine: Life and death
				\item AI: Criminal justice, finance, healthcare
			\end{itemize}
			\item \textbf{Rapid Technological Change:}
			\begin{itemize}
				\item Medicine: Mechanical ventilation, transplants
				\item AI: Rapid advancement, unpredictable effects
			\end{itemize}
			\item \textbf{Need for Oversight:}
			\begin{itemize}
				\item Medicine: Ethics committees
				\item AI: Governance structures needed
			\end{itemize}
			\item \textbf{Professional Responsibility:}
			\begin{itemize}
				\item Medicine: Doctors' oaths
				\item AI: Developers' responsibilities
			\end{itemize}
		\end{enumerate}
		
		\vspace{0.1cm}
		\textbf{Key Difference:}
		\begin{itemize}
			\item Medical ethics has centuries of development
			\item AI ethics is still evolving
			\item Can learn from medical ethics
		\end{itemize}
		
		\textcolor{myblue}{\textbf{Key Insight:}} AI ethics can \textcolor{myblue}{learn from established medical ethics} structures!
	\end{frame}
	
	% ============================================================
	% SECTION 5: PRINCIPLES OF AI ETHICS - SLIDES 41-55
	% ============================================================
	
	\begin{frame}{Slide 41: AI Ethics Principles Overview}
		\fontsize{6.5}{7.5}\selectfont
		\textbf{Five Core Principles of AI Ethics:}
		
		\vspace{0.1cm}
		\textbf{The Five Principles:}
		\begin{enumerate}
			\item \textbf{Nonmaleficence:}
			\begin{itemize}
				\item Obligation to avoid harming others
				\item "Do no harm"
				\item Prohibition of unjustified risks
			\end{itemize}
			\item \textbf{Beneficence:}
			\begin{itemize}
				\item Obligation to act for the benefit of others
				\item "Do good"
				\item Active help and support
			\end{itemize}
			\item \textbf{Justice:}
			\begin{itemize}
				\item Fairness, equality, impartiality
				\item Equitable distribution
				\item Non-discrimination
			\end{itemize}
			\item \textbf{Autonomy:}
			\begin{itemize}
				\item Respect for informed, un-coerced decisions
				\item Individual control
				\item Non-manipulation
			\end{itemize}
			\item \textbf{Explainability:}
			\begin{itemize}
				\item Transparent and understandable decisions
				\item Accountability
				\item Trust-building
			\end{itemize}
		\end{enumerate}
		
		\textcolor{myblue}{\textbf{Key Insight:}} Five principles = \textcolor{myblue}{nonmaleficence, beneficence, justice, autonomy, explainability}!
	\end{frame}
	
	\begin{frame}{Slide 42: Nonmaleficence in Detail}
		\fontsize{6.5}{7.5}\selectfont
		\textbf{The Principle of Nonmaleficence:}
		
		\vspace{0.1cm}
		\textbf{Core Meaning:}
		\begin{itemize}
			\item Obligation to avoid harming others
			\item Prohibition on imposing unjustified risks
			\item "First, do no harm" (primum non nocere)
		\end{itemize}
		
		\vspace{0.1cm}
		\textbf{Key Aspects:}
		\begin{enumerate}
			\item \textbf{Avoid Actual Harm:}
			\begin{itemize}
				\item Don't directly harm others
				\item Example: Algorithms that discriminate
				\item Actions that cause injury
			\end{itemize}
			\item \textbf{Avoid Unjustified Risk:}
			\begin{itemize}
				\item Don't impose unnecessary risks
				\item Example: Untested algorithms
				\item Risks that outweigh benefits
			\end{itemize}
			\item \textbf{Proportionality:}
			\begin{itemize}
				\item Some risk is permissible
				\item If benefits justify risk
				\item If alternatives unavailable
			\end{itemize}
		\end{enumerate}
		
		\textcolor{myblue}{\textbf{Key Insight:}} Nonmaleficence = \textcolor{myblue}{"do no harm"} - avoid harming others!
	\end{frame}
	
	\begin{frame}{Slide 43: Risk Justification}
		\fontsize{6.5}{7.5}\selectfont
		\textbf{When is Risk Acceptable?}
		
		\vspace{0.1cm}
		\textbf{Principles of Risk Acceptance:}
		\begin{enumerate}
			\item \textbf{Benefit Justification:}
			\begin{itemize}
				\item Risk is acceptable if benefits justify it
				\item Example: Clinical trials
				\item Risk is part of medical progress
			\end{itemize}
			\item \textbf{No Alternatives:}
			\begin{itemize}
				\item Risk acceptable if no alternatives
				\item Example: Emergency procedures
				\item Only option available
			\end{itemize}
			\item \textbf{Who Bears the Risk:}
			\begin{itemize}
				\item More acceptable if those at risk benefit
				\item Example: Clinical trial participants
				\item Unacceptable if others bear risk
			\end{itemize}
		\end{enumerate}
		
		\vspace{0.1cm}
		\textbf{AI Application:}
		\begin{itemize}
			\item Imposing algorithmic risk on people who don't benefit
			\item Unacceptable to experiment on populations
			\item Especially if they don't gain
		\end{itemize}
		
		\textcolor{myblue}{\textbf{Key Insight:}} Risk acceptable if \textcolor{myblue}{benefits justify it and risk-bearers benefit}!
	\end{frame}
	
	\begin{frame}{Slide 44: Beneficence in Detail}
		\fontsize{6.5}{7.5}\selectfont
		\textbf{The Principle of Beneficence:}
		
		\vspace{0.1cm}
		\textbf{Core Meaning:}
		\begin{itemize}
			\item Moral obligation to act for the benefit of others
			\item Requires taking positive steps to help
			\item Moves beyond simply refraining from harm
		\end{itemize}
		
		\vspace{0.1cm}
		\textbf{Key Aspects:}
		\begin{enumerate}
			\item \textbf{Active Help:}
			\begin{itemize}
				\item Donating to charity
				\item Volunteering
				\item Providing resources and assistance
			\end{itemize}
			\item \textbf{Beyond Nonmaleficence:}
			\begin{itemize}
				\item Not just "do no harm"
				\item "Do good"
				\item Active beneficence
			\end{itemize}
			\item \textbf{Deontological Duty:}
			\begin{itemize}
				\item Duty-based obligation
				\item Not just about maximizing utility
				\item Required beyond consequences
			\end{itemize}
		\end{enumerate}
		
		\textcolor{myblue}{\textbf{Key Insight:}} Beneficence = \textcolor{myblue}{"do good"} - actively help others!
	\end{frame}
	
	\begin{frame}{Slide 45: Beneficence Limitations}
		\fontsize{6.5}{7.5}\selectfont
		\textbf{Limits to the Duty of Beneficence:}
		
		\vspace{0.1cm}
		\textbf{Reasonable Constraints:}
		\begin{enumerate}
			\item \textbf{Scarce Resources:}
			\begin{itemize}
				\item Cannot help everyone
				\item Limited resources
				\item Prioritization needed
			\end{itemize}
			\item \textbf{Competing Obligations:}
			\begin{itemize}
				\item Multiple duties to balance
				\item Professional responsibilities
				\item Personal obligations
			\end{itemize}
			\item \textbf{Reasonableness:}
			\begin{itemize}
				\item There's only so much ethics can demand
				\item Demandingness constraint
				\item Proportionality matters
			\end{itemize}
			\item \textbf{Recipient Autonomy:}
			\begin{itemize}
				\item Unwanted "benefits" may be rejected
				\item Right to decide what's best
				\item Respect individual choices
			\end{itemize}
		\end{enumerate}
		
		\textcolor{myblue}{\textbf{Key Insight:}} Beneficence has \textcolor{myblue}{reasonable limits} - cannot help everyone!
	\end{frame}
	
	\begin{frame}{Slide 46: Justice in AI}
		\fontsize{6.5}{7.5}\selectfont
		\textbf{Justice in the Context of AI:}
		
		\vspace{0.1cm}
		\textbf{Core Meaning:}
		\begin{itemize}
			\item Fairness, equality, impartiality, proportionality
			\item Giving each person their proper due
			\item Upholding duties toward fairness and equality
		\end{itemize}
		
		\vspace{0.1cm}
		\textbf{Key Applications:}
		\begin{enumerate}
			\item \textbf{Fair Treatment:}
			\begin{itemize}
				\item No discrimination based on protected characteristics
				\item Equal opportunities
				\item Impartial algorithms
			\end{itemize}
			\item \textbf{Equitable Distribution:}
			\begin{itemize}
				\item Benefits and burdens fairly allocated
				\item Access to AI benefits
				\item Protection from AI harms
			\end{itemize}
			\item \textbf{Accountability:}
			\begin{itemize}
				\item Those responsible should be accountable
				\item Just processes
				\item Fair consequences
			\end{itemize}
		\end{enumerate}
		
		\textcolor{myblue}{\textbf{Key Insight:}} Justice = \textcolor{myblue}{fairness, equality, impartiality, and proportionality}!
	\end{frame}
	
	\begin{frame}{Slide 47: Types of Justice}
		\fontsize{6.5}{7.5}\selectfont
		\textbf{Four Types of Justice:}
		
		\vspace{0.1cm}
		\textbf{1. Procedural Justice:}
		\begin{itemize}
			\item Fair processes and impartiality
			\item Right structures in place
			\item Rule of law
		\end{itemize}
		
		\vspace{0.1cm}
		\textbf{2. Distributive Justice:}
		\begin{itemize}
			\item Equitable allocation of benefits and burdens
			\item Fair distribution of resources
			\item Access to opportunities
		\end{itemize}
		
		\vspace{0.1cm}
		\textbf{3. Restorative Justice:}
		\begin{itemize}
			\item Repairing harms
			\item Reconciling victims and offenders
			\item Making amends
		\end{itemize}
		
		\vspace{0.1cm}
		\textbf{4. Interactional Justice:}
		\begin{itemize}
			\item Respect and fairness between individuals
			\item Interpersonal treatment
			\item Dignity and respect
		\end{itemize}
		
		\textcolor{myblue}{\textbf{Key Insight:}} Four types of justice = \textcolor{myblue}{procedural, distributive, restorative, interactional}!
	\end{frame}
	
	\begin{frame}{Slide 48: Autonomy in Detail}
		\fontsize{6.5}{7.5}\selectfont
		\textbf{The Principle of Autonomy:}
		
		\vspace{0.1cm}
		\textbf{Core Meaning:}
		\begin{itemize}
			\item Capacity to make informed, un-coerced decisions
			\item Respecting others' abilities to determine their own course
			\item Non-manipulation
		\end{itemize}
		
		\vspace{0.1cm}
		\textbf{Key Aspects:}
		\begin{enumerate}
			\item \textbf{Informed Decisions:}
			\begin{itemize}
				\item People need adequate information
				\item Transparency about options
				\item Understanding consequences
			\end{itemize}
			\item \textbf{Un-coerced Decisions:}
			\begin{itemize}
				\item No coercion or manipulation
				\item Free choice
				\item Voluntary consent
			\end{itemize}
			\item \textbf{Capacity:}
			\begin{itemize}
				\item Must have capacity to decide
				\item Exceptions: children, unconscious patients
				\item Surrogate decision makers when needed
			\end{itemize}
		\end{enumerate}
		
		\textcolor{myblue}{\textbf{Key Insight:}} Autonomy = \textcolor{myblue}{informed, un-coerced decisions}!
	\end{frame}
	
	\begin{frame}{Slide 49: Autonomy in Healthcare}
		\fontsize{6.5}{7.5}\selectfont
		\textbf{Autonomy in Medical Context:}
		
		\vspace{0.1cm}
		\textbf{Informed Consent:}
		\begin{itemize}
			\item Patients have right to voluntary informed consent
			\item Right to refuse treatment
			\item Doctor's duty to inform
			\item Patient's right to decide
		\end{itemize}
		
		\vspace{0.1cm}
		\textbf{What Undermines Autonomy:}
		\begin{enumerate}
			\item \textbf{Coercion:}
			\begin{itemize}
				\item Forcing decisions
				\item Threats or pressure
				\item Removing choice
			\end{itemize}
			\item \textbf{Deception:}
			\begin{itemize}
				\item Lying or misleading
				\item Withholding information
				\item Manipulating understanding
			\end{itemize}
			\item \textbf{Undue Influence:}
			\begin{itemize}
				\item Excessive persuasion
				\item Exploitation
				\item Imbalance of power
			\end{itemize}
		\end{enumerate}
		
		\textcolor{myblue}{\textbf{Key Insight:}} Autonomy requires \textcolor{myblue}{informed consent and no coercion}!
	\end{frame}
	
	\begin{frame}{Slide 50: Autonomy in AI}
		\fontsize{6.5}{7.5}\selectfont
		\textbf{Autonomy in the Context of AI:}
		
		\vspace{0.1cm}
		\textbf{Ethical AI and Autonomy:}
		\begin{itemize}
			\item Ethical AI respects people's autonomy
			\item No coercive or manipulative tactics
			\item Technology should help people further their own goals
			\item Not serve third-party interests
		\end{itemize}
		
		\vspace{0.1cm}
		\textbf{Risks to Autonomy:}
		\begin{enumerate}
			\item \textbf{Manipulation:}
			\begin{itemize}
				\item Recommendation algorithms
				\item Targeted content
				\item Behavioral influence
			\end{itemize}
			\item \textbf{Coercion:}
			\begin{itemize}
				\item Forced compliance
				\item No alternative
				\item System lock-in
			\end{itemize}
			\item \textbf{Lack of Transparency:}
			\begin{itemize}
				\item Users don't understand decisions
				\item Cannot make informed choices
				\item Autonomy undermined
			\end{itemize}
		\end{enumerate}
		
		\textcolor{myblue}{\textbf{Key Insight:}} Ethical AI respects \textcolor{myblue}{autonomy by avoiding coercion and manipulation}!
	\end{frame}
	
	\begin{frame}{Slide 51: Explainability in AI}
		\fontsize{6.5}{7.5}\selectfont
		\textbf{The Principle of Explainability:}
		
		\vspace{0.1cm}
		\textbf{Core Meaning:}
		\begin{itemize}
			\item AI systems should provide transparent explanations
			\item Understandable reasons for actions or decisions
			\item Also called "explicability"
		\end{itemize}
		
		\vspace{0.1cm}
		\textbf{Why Explainability Matters:}
		\begin{enumerate}
			\item \textbf{Accountability:}
			\begin{itemize}
				\item Hold designers and companies accountable
				\item Identify errors and biases
				\item Rectify unfair practices
			\end{itemize}
			\item \textbf{Trust:}
			\begin{itemize}
				\item Trust is fundamental for adoption
				\item Users trust understandable systems
				\item Transparency builds confidence
			\end{itemize}
			\item \textbf{Ethical Compliance:}
			\begin{itemize}
				\item Ensure adherence to ethical guidelines
				\item Respect individual rights
				\item Regulatory compliance
			\end{itemize}
		\end{enumerate}
		
		\textcolor{myblue}{\textbf{Key Insight:}} Explainability = \textcolor{myblue}{transparent and understandable explanations}!
	\end{frame}
	
	\begin{frame}{Slide 52: Types of Explainability}
		\fontsize{6.5}{7.5}\selectfont
		\textbf{Levels and Types of Explainability:}
		
		\vspace{0.1cm}
		\textbf{1. Local Explainability:}
		\begin{itemize}
			\item Explains decisions on a single instance
			\item "Why was THIS decision made?"
			\item Example: Why was THIS loan denied?
			\item Specific to one prediction
		\end{itemize}
		
		\vspace{0.1cm}
		\textbf{2. Global Explainability:}
		\begin{itemize}
			\item Explains overall model behavior
			\item "How does the model work generally?"
			\item Comprehensive understanding
			\item Across all inputs
		\end{itemize}
		
		\vspace{0.1cm}
		\textbf{3. Model-Specific Explainability:}
		\begin{itemize}
			\item Techniques tailored to architecture
			\item Example: Decision tree rules
			\item Neural networks require different methods
		\end{itemize}
		
		\vspace{0.1cm}
		\textbf{4. Model-Agnostic Explainability:}
		\begin{itemize}
			\item Works with any model
			\item Versatile and general
			\item LIME, SHAP are examples
		\end{itemize}
		
		\textcolor{myblue}{\textbf{Key Insight:}} Explainability has \textcolor{myblue}{different levels and types} for different needs!
	\end{frame}
	
	\begin{frame}{Slide 53: Counterfactual Explanations}
		\fontsize{6.5}{7.5}\selectfont
		\textbf{Counterfactual Explanations:}
		
		\vspace{0.1cm}
		\textbf{What They Are:}
		\begin{itemize}
			\item Hypothetical scenarios contrasting with actual decisions
			\item Show conditions for different outcomes
			\item Provide actionable insights
		\end{itemize}
		
		\vspace{0.1cm}
		\textbf{Example - Loan Denial:}
		\begin{itemize}
			\item Actual: Loan denied
			\item Counterfactual: "You would have been approved if:
			\begin{itemize}
				\item You had \$10,000 more in the bank
				\item OR you earned \$1,000 more per month
			\end{itemize}
		\end{itemize}
		
		\vspace{0.1cm}
		\textbf{Benefits:}
		\begin{enumerate}
			\item \textbf{Transparency:}
			\begin{itemize}
				\item Users understand decision factors
				\item See what influences outcomes
			\end{itemize}
			\item \textbf{Actionable Insights:}
			\begin{itemize}
				\item Know what to change
				\item Path to different outcome
				\item Empowerment
			\end{itemize}
			\item \textbf{Fairness:}
			\begin{itemize}
				\item Detect biases
				\item Understand disparate treatment
				\item Identify improvement areas
			\end{itemize}
		\end{enumerate}
		
		\textcolor{myblue}{\textbf{Key Insight:}} Counterfactual explanations show \textcolor{myblue}{what would lead to different outcomes}!
	\end{frame}
	
	\begin{frame}{Slide 54: Explainability Challenges}
		\fontsize{6.5}{7.5}\selectfont
		\textbf{Challenges with Explainability:}
		
		\vspace{0.1cm}
		\textbf{Key Challenges:}
		\begin{enumerate}
			\item \textbf{What Counts as Explanation?}
			\begin{itemize}
				\item No consensus on definition
				\item Different needs for different audiences
				\item Experts vs. general public
			\end{itemize}
			\item \textbf{Who Gets Explanation?}
			\begin{itemize}
				\item Affected individuals
				\item Regulators
				\item General public
				\item Different levels needed
			\end{itemize}
			\item \textbf{Trade-offs:}
			\begin{itemize}
				\item Accuracy vs. explainability
				\item Simple models are more explainable
				\item Complex models are more accurate
			\end{itemize}
			\item \textbf{Technical Difficulty:}
			\begin{itemize}
				\item Neural networks are "black boxes"
				\item Hard to trace decisions
				\item Technical solutions needed
			\end{itemize}
		\end{enumerate}
		
		\textcolor{myblue}{\textbf{Key Insight:}} Explainability faces \textcolor{myblue}{definitional, audience, and technical challenges}!
	\end{frame}
	
	\begin{frame}{Slide 55: Principles Summary}
		\fontsize{6.5}{7.5}\selectfont
		\textbf{Summary of AI Ethics Principles:}
		
		\vspace{0.1cm}
		\textbf{Five Core Principles:}
		\begin{enumerate}
			\item \textbf{Nonmaleficence:}
			\begin{itemize}
				\item "Do no harm"
				\item Avoid unjustified risks
				\item Proportional risk-benefit
			\end{itemize}
			\item \textbf{Beneficence:}
			\begin{itemize}
				\item "Do good"
				\item Active help
				\item Advance welfare
			\end{itemize}
			\item \textbf{Justice:}
			\begin{itemize}
				\item Fairness and equality
				\item Impartiality
				\item Non-discrimination
			\end{itemize}
			\item \textbf{Autonomy:}
			\begin{itemize}
				\item Informed decisions
				\item Non-coercion
				\item Respect for persons
			\end{itemize}
			\item \textbf{Explainability:}
			\begin{itemize}
				\item Transparency
				\item Understandable decisions
				\item Accountability
			\end{itemize}
		\end{enumerate}
		
		\textcolor{myblue}{\textbf{Key Insight:}} Five principles provide \textcolor{myblue}{comprehensive ethical framework} for AI!
	\end{frame}
	
	% ============================================================
	% SECTION 6: BIAS, DISCRIMINATION, AND FAIRNESS - SLIDES 56-75
	% ============================================================
	
	\begin{frame}{Slide 56: Bias in AI}
		\fontsize{6.5}{7.5}\selectfont
		\textbf{Understanding Bias in AI Systems:}
		
		\vspace{0.1cm}
		\textbf{Definition:}
		\begin{itemize}
			\item Systemic deviation in output or impact
			\item Compared to a desired norm or standard
			\item Algorithm is biased when it deviates from intended function
		\end{itemize}
		
		\vspace{0.1cm}
		\textbf{Example - Hiring Algorithm:}
		\begin{itemize}
			\item Designed to identify best candidates
			\item But recommends based on race or sex
			\item Deviates from intended purpose
			\item This is problematic bias
		\end{itemize}
		
		\vspace{0.1cm}
		\textbf{Key Insight:}
		\begin{itemize}
			\item No algorithm is perfectly objective
			\item Every algorithm reflects embedded values
			\item Values shape what's considered valuable
		\end{itemize}
		
		\textcolor{myblue}{\textbf{Key Insight:}} AI bias = \textcolor{myblue}{systemic deviation from intended function}!
	\end{frame}
	
	\begin{frame}{Slide 57: Justifiable vs Problematic Bias}
		\fontsize{6.5}{7.5}\selectfont
		\textbf{When is Bias Acceptable?}
		
		\vspace{0.1cm}
		\textbf{Justifiable Bias:}
		\begin{itemize}
			\item Part of well-designed AI
			\item Optimizes for relevant criteria
			\item Example: Loan eligibility prioritizing bank balance
			\item Relevant to loan repayment
		\end{itemize}
		
		\vspace{0.1cm}
		\textbf{Problematic Bias:}
		\begin{itemize}
			\item Results in unfairness
			\item Disadvantages for unjustifiable reasons
			\item Based on irrelevant characteristics
			\item Example: Race or sex discrimination
		\end{itemize}
		
		\vspace{0.1cm}
		\textbf{Statistically Unbiased not equal to Ethically Acceptable:}
		\begin{itemize}
			\item Even unbiased algorithms can cause harm
			\item Example: Charging exorbitant fees to everyone
			\item Statistically fair but ethically problematic
		\end{itemize}
		
		\textcolor{myblue}{\textbf{Key Insight:}} Not all biases are \textcolor{myblue}{problematic} - some are justifiable!
	\end{frame}
	
	\begin{frame}{Slide 58: Sources of Bias - Overview}
		\fontsize{6.5}{7.5}\selectfont
		\textbf{Four Primary Sources of Algorithmic Bias:}
		
		\vspace{0.1cm}
		\textbf{1. Problem Specification:}
		\begin{itemize}
			\item Goals contain inherent problems
			\item Target variables miss real-world objectives
			\item Example: Healthcare expenditure proxy
		\end{itemize}
		
		\vspace{0.1cm}
		\textbf{2. Data:}
		\begin{itemize}
			\item Historical bias
			\item Sampling bias
			\item Selective labels problem
		\end{itemize}
		
		\vspace{0.1cm}
		\textbf{3. Modeling, Validation, and Design:}
		\begin{itemize}
			\item Optimization choices
			\item Model selection
			\item Subgroup consideration
		\end{itemize}
		
		\vspace{0.1cm}
		\textbf{4. Deployment:}
		\begin{itemize}
			\item Human deference
			\item Unintended effects
			\item Context changes
		\end{itemize}
		
		\textcolor{myblue}{\textbf{Key Insight:}} Bias can enter at \textcolor{myblue}{any stage} of AI development!
	\end{frame}
	
	\begin{frame}{Slide 59: Problem Specification Bias}
		\fontsize{6.5}{7.5}\selectfont
		\textbf{Bias in Problem Specification:}
		
		\vspace{0.1cm}
		\textbf{How It Occurs:}
		\begin{enumerate}
			\item \textbf{Complex Goal Operationalization:}
			\begin{itemize}
				\item Hard to capture real-world objectives
				\item Selected target variables may be poor proxies
				\item Example: Bank eliminating all risk
			\end{itemize}
			\item \textbf{Unintended Consequences:}
			\begin{itemize}
				\item Bank lends only to affluent
				\item Excludes others
				\item Fairness issue
			\end{itemize}
			\item \textbf{Proxy Problems:}
			\begin{itemize}
				\item Healthcare expenditure proxy
				\item Favored richest patients, not sickest
				\item Systematically biased
			\end{itemize}
		\end{enumerate}
		
		\vspace{0.1cm}
		\textbf{Key Lesson:}
		\begin{itemize}
			\item Careful problem definition is essential
			\item Proxies must be thoughtfully chosen
			\item Goals should reflect true objectives
		\end{itemize}
		
		\textcolor{myblue}{\textbf{Key Insight:}} Problem specification can introduce \textcolor{myblue}{bias through poor operationalization}!
	\end{frame}
	
	\begin{frame}{Slide 60: Data-Related Biases}
		\fontsize{6.5}{7.5}\selectfont
		\textbf{Biases in Data:}
		
		\vspace{0.1cm}
		\textbf{Historical Bias:}
		\begin{itemize}
			\item Data reflects past discrimination
			\item Example: Women excluded from banking
			\item Algorithm favors men
			\item Perpetuates historical patterns
		\end{itemize}
		
		\vspace{0.1cm}
		\textbf{Selective Labels Problem:}
		\begin{itemize}
			\item Missing counterfactual outcomes
			\item No data on rejected applicants
			\item Can't know if better candidates existed
			\item Perpetuates existing patterns
		\end{itemize}
		
		\vspace{0.1cm}
		\textbf{Sampling Bias:}
		\begin{itemize}
			\item Non-random data sample
			\item Trends don't generalize
			\item Example: Data from young men only
			\item May not apply to women or elderly
		\end{itemize}
		
		\textcolor{myblue}{\textbf{Key Insight:}} Data biases include \textcolor{myblue}{historical, selective labels, and sampling bias}!
	\end{frame}
	
	\begin{frame}{Slide 61: Modeling and Design Bias}
		\fontsize{6.5}{7.5}\selectfont
		\textbf{Bias in Modeling, Validation, and Design:}
		
		\vspace{0.1cm}
		\textbf{How Bias Emerges:}
		\begin{enumerate}
			\item \textbf{Optimization Functions:}
			\begin{itemize}
				\item Choice of objective
				\item What is being optimized?
				\item May favor certain outcomes
			\end{itemize}
			\item \textbf{Regression Models:}
			\begin{itemize}
				\item Different models have different biases
				\item Choice matters
				\item Subgroup effects
			\end{itemize}
			\item \textbf{Subgroup Consideration:}
			\begin{itemize}
				\item How are subgroups handled?
				\item Are differences recognized?
				\item Fairness across groups
			\end{itemize}
			\item \textbf{Information Presentation:}
			\begin{itemize}
				\item How is data presented?
				\item Visual biases
				\item Interface effects
			\end{itemize}
		\end{enumerate}
		
		\textcolor{myblue}{\textbf{Key Insight:}} Modeling choices can \textcolor{myblue}{introduce bias} even with good data!
	\end{frame}
	
	\begin{frame}{Slide 62: Search Engine Bias Example}
		\fontsize{6.5}{7.5}\selectfont
		\textbf{Search Engine Bias Example:}
		
		\vspace{0.1cm}
		\textbf{The Scenario:}
		\begin{itemize}
			\item Search engine for financial products
			\item Designed to help users find best options
			\item Favors popular items
			\item Popularity leads to more exposure
		\end{itemize}
		
		\vspace{0.1cm}
		\textbf{The Bias:}
		\begin{itemize}
			\item Perpetuates popularity cycles
			\item Quality may not matter
			\item Less popular but better options hidden
			\item Self-reinforcing bias
		\end{itemize}
		
		\vspace{0.1cm}
		\textbf{Lesson:}
		\begin{itemize}
			\item Popularity not equal to quality
			\item Algorithmic bias can be subtle
			\item Design choices matter
			\item Fairness implications
		\end{itemize}
		
		\textcolor{myblue}{\textbf{Key Insight:}} Design choices can \textcolor{myblue}{create self-reinforcing biases}!
	\end{frame}
	
	\begin{frame}{Slide 63: Deployment Bias}
		\fontsize{6.5}{7.5}\selectfont
		\textbf{Bias in Deployment:}
		
		\vspace{0.1cm}
		\textbf{How It Occurs:}
		\begin{enumerate}
			\item \textbf{Human Deference:}
			\begin{itemize}
				\item People defer to algorithm suggestions
				\item Even when limitations known
				\item Convenience factor
			\end{itemize}
			\item \textbf{Reasons for Deference:}
			\begin{itemize}
				\item Time-saving
				\item Perceived objectivity
				\item Self-doubt
				\item Responsibility shielding
			\end{itemize}
			\item \textbf{Unintended Effects:}
			\begin{itemize}
				\item Individuals relinquish responsibility
				\item Blame shifting
				\item Behavior changes
			\end{itemize}
		\end{enumerate}
		
		\vspace{0.1cm}
		\textbf{Key Lesson:}
		\begin{itemize}
			\item Oversight is essential
			\item Human judgment remains important
			\item Don't blindly trust automation
		\end{itemize}
		
		\textcolor{myblue}{\textbf{Key Insight:}} Deployment can introduce \textcolor{myblue}{bias through human deference} to algorithms!
	\end{frame}
	
	\begin{frame}{Slide 64: Automation Deference}
		\fontsize{6.5}{7.5}\selectfont
		\textbf{Why People Defer to Automation:}
		
		\vspace{0.1cm}
		\textbf{Convenience:}
		\begin{itemize}
			\item Time-saving
			\item Effort reduction
			\item Easy to follow
		\end{itemize}
		
		\vspace{0.1cm}
		\textbf{Perceived Objectivity:}
		\begin{itemize}
			\item Algorithms appear more objective
			\item People know their own fallibilities
			\item Creates self-doubt
		\end{itemize}
		
		\vspace{0.1cm}
		\textbf{Responsibility Shielding:}
		\begin{itemize}
			\item Deferring shares responsibility
			\item Going against algorithm exposes people
			\item Blame avoidance
		\end{itemize}
		
		\vspace{0.1cm}
		\textbf{Consequences:}
		\begin{itemize}
			\item Unintended effects
			\item Reduced oversight
			\item Bias perpetuation
			\item Accountability gaps
		\end{itemize}
		
		\textcolor{myblue}{\textbf{Key Insight:}} People defer to algorithms for \textcolor{myblue}{convenience, objectivity perception, and responsibility shielding}!
	\end{frame}
	
	\begin{frame}{Slide 65: Discrimination in Law}
		\fontsize{6.5}{7.5}\selectfont
		\textbf{When Bias Becomes Discrimination:}
		
		\vspace{0.1cm}
		\textbf{Legal Perspective:}
		\begin{itemize}
			\item Depends on jurisdiction
			\item When it disfavors based on protected characteristics
			\item Race, sex, ethnicity, age
			\item Other classifications protected by law
		\end{itemize}
		
		\vspace{0.1cm}
		\textbf{Protected Characteristics:}
		\begin{enumerate}
			\item Race
			\item Sex/Gender
			\item Ethnicity
			\item Age
			\item Religion
			\item Disability
			\item Sexual orientation
		\end{enumerate}
		
		\vspace{0.1cm}
		\textbf{Key Point:}
		\begin{itemize}
			\item Such disadvantages typically violate legal protections
			\item Proactive measures needed
			\item Continuous monitoring
		\end{itemize}
		
		\textcolor{myblue}{\textbf{Key Insight:}} Discrimination = \textcolor{myblue}{disfavoring based on protected characteristics} in violation of law!
	\end{frame}
	
	\begin{frame}{Slide 66: Fairness in AI}
		\fontsize{6.5}{7.5}\selectfont
		\textbf{Understanding Fairness in AI:}
		
		\vspace{0.1cm}
		\textbf{Definition:}
		\begin{itemize}
			\item Absence of bias or preference
			\item Based on irrelevant characteristics
			\item Algorithm is fair when it doesn't exhibit problematic biases
		\end{itemize}
		
		\vspace{0.1cm}
		\textbf{Two Types:}
		\begin{enumerate}
			\item \textbf{Group Fairness:}
			\begin{itemize}
				\item Statistical criteria across groups
				\item Demographics mirror population
				\item Example: 51% female = 51% loan recipients
			\end{itemize}
			\item \textbf{Individual Fairness:}
			\begin{itemize}
				\item Treating similar individuals similarly
				\item Challenge: Defining "similar"
				\item Focus on individual treatment
			\end{itemize}
		\end{enumerate}
		
		\textcolor{myblue}{\textbf{Key Insight:}} Fairness = \textcolor{myblue}{absence of problematic biases}!
	\end{frame}
	
	\begin{frame}{Slide 67: Fairness Impossibility}
		\fontsize{6.5}{7.5}\selectfont
		\textbf{Why Fairness Can't Be Automated:}
		
		\vspace{0.1cm}
		\textbf{Mathematical Impossibility:}
		\begin{itemize}
			\item When base rates are unequal
			\item Cannot satisfy all fairness measures simultaneously
			\item Demographic parity, predictive rate parity, equal opportunity
			\item Trade-offs are unavoidable
		\end{itemize}
		
		\vspace{0.1cm}
		\textbf{Example - Criminal Risk Assessment:}
		\begin{itemize}
			\item Men statistically more likely to commit certain crimes
			\item Algorithm assigns higher risk to man
			\item Man argues unfair treatment
			\item Woman argues unfair if assigned equal score
		\end{itemize}
		
		\vspace{0.1cm}
		\textbf{Key Point:}
		\begin{itemize}
			\item Fairness is a moral judgment
			\item Not purely mathematical
			\item Compromises and trade-offs needed
		\end{itemize}
		
		\textcolor{myblue}{\textbf{Key Insight:}} Fairness cannot be automated when \textcolor{myblue}{base rates are unequal}!
	\end{frame}
	
	\begin{frame}{Slide 68: Procedural Fairness}
		\fontsize{6.5}{7.5}\selectfont
		\textbf{Procedural Fairness in AI:}
		
		\vspace{0.1cm}
		\textbf{Definition:}
		\begin{itemize}
			\item Fair processes and impartial procedures
			\item Not just fair outcomes
			\item Right structures in place
		\end{itemize}
		
		\vspace{0.1cm}
		\textbf{JUSTICE System Analogy:}
		\begin{itemize}
			\item Procedural fairness: rule of law
			\item Outcome fairness: appropriate punishments
			\item Mistakes can be justifiable with fair processes
			\item Ways to right wrongs
		\end{itemize}
		
		\vspace{0.1cm}
		\textbf{AI Application:}
		\begin{itemize}
			\item Ethics committees
			\item Established procedures
			\item Appeal processes
			\item Accountability structures
		\end{itemize}
		
		\textcolor{myblue}{\textbf{Key Insight:}} Procedural fairness = \textcolor{myblue}{fair processes and impartial procedures}!
	\end{frame}
	
	\begin{frame}{Slide 69: Avoiding Problematic Biases}
		\fontsize{6.5}{7.5}\selectfont
		\textbf{Strategies to Avoid Problematic Biases:}
		
		\vspace{0.1cm}
		\textbf{Key Approaches:}
		\begin{enumerate}
			\item \textbf{Technological Solutions:}
			\begin{itemize}
				\item Fairness toolkits (Aequitas, AIF360)
				\item Internal auditing systems
				\item Bias reports from providers
				\item But technology alone is insufficient
			\end{itemize}
			\item \textbf{Data Quality:}
			\begin{itemize}
				\item Diverse, updated, accurate data
				\item Representative samples
				\item Free from past discriminatory tendencies
			\end{itemize}
			\item \textbf{Auditing:}
			\begin{itemize}
				\item Algorithmic auditing is best approach
				\item Private companies offer services
				\item Independent assessment
			\end{itemize}
		\end{enumerate}
		
		\textcolor{myblue}{\textbf{Key Insight:}} Bias mitigation requires \textcolor{myblue}{multiple approaches beyond technology}!
	\end{frame}
	
	\begin{frame}{Slide 70: Data Quality Importance}
		\fontsize{6.5}{7.5}\selectfont
		\textbf{Data Quality and Bias Prevention:}
		
		\vspace{0.1cm}
		\textbf{Key Data Quality Factors:}
		\begin{enumerate}
			\item \textbf{Diversity:}
			\begin{itemize}
				\item Include all relevant groups
				\item Avoid homogeneity
				\item Represent all populations
			\end{itemize}
			\item \textbf{Accuracy:}
			\begin{itemize}
				\item Correct and reliable data
				\item Avoid errors
				\item Verify sources
			\end{itemize}
			\item \textbf{Representativeness:}
			\begin{itemize}
				\item Reflect target population
				\item Avoid sampling bias
				\item Generalizable results
			\end{itemize}
			\item \textbf{Currency:}
			\begin{itemize}
				\item Up-to-date data
				\item Reflect current reality
				\item Avoid outdated patterns
			\end{itemize}
		\end{enumerate}
		
		\vspace{0.1cm}
		\textbf{Synthetic Data:}
		\begin{itemize}
			\item Fabricated datasets
			\item No privacy risks
			\item Can be designed to avoid biases
			\item May be less precise
		\end{itemize}
		
		\textcolor{myblue}{\textbf{Key Insight:}} Data quality helps \textcolor{myblue}{avoid biases} through diversity and representativeness!
	\end{frame}
	
	\begin{frame}{Slide 71: Algorithmic Auditing}
		\fontsize{6.5}{7.5}\selectfont
		\textbf{The Importance of Algorithmic Auditing:}
		
		\vspace{0.1cm}
		\textbf{Why Auditing Matters:}
		\begin{itemize}
			\item Likely the best way to identify biases
			\item Fresh and diverse perspective
			\item Independent assessment
			\item Continuous improvement
		\end{itemize}
		
		\vspace{0.1cm}
		\textbf{What Auditing Includes:}
		\begin{enumerate}
			\item \textbf{Technological Tools:}
			\begin{itemize}
				\item Statistical analyses
				\item Fairness metrics
				\item Bias detection
			\end{itemize}
			\item \textbf{Human Review:}
			\begin{itemize}
				\item Diverse perspectives
				\item Contextual understanding
				\item Qualitative assessment
			\end{itemize}
			\item \textbf{Documentation:}
			\begin{itemize}
				\item Findings and recommendations
				\item Remediation plans
				\item Continuous monitoring
			\end{itemize}
		\end{enumerate}
		
		\textcolor{myblue}{\textbf{Key Insight:}} Algorithmic auditing is \textcolor{myblue}{best way to identify and correct biases}!
	\end{frame}
	
	\begin{frame}{Slide 72: Ethical Committees}
		\fontsize{6.5}{7.5}\selectfont
		\textbf{The Role of Ethical Committees:}
		
		\vspace{0.1cm}
		\textbf{Purpose:}
		\begin{itemize}
			\item Discuss possible problematic biases
			\item Make decisions about trade-offs
			\item Consider multiple ethical frameworks
			\item Ensure procedural fairness
		\end{itemize}
		
		\vspace{0.1cm}
		\textbf{Benefits:}
		\begin{enumerate}
			\item \textbf{Diverse Perspectives:}
			\begin{itemize}
				\item Different viewpoints
				\item Multiple expertise areas
				\item Balanced considerations
			\end{itemize}
			\item \textbf{Structured Decision-Making:}
			\begin{itemize}
				\item Systematic review
				\item Transparent processes
				\item Documented reasoning
			\end{itemize}
			\item \textbf{Accountability:}
			\begin{itemize}
				\item Clear responsibility
				\item Audit trail
				\item Defensible decisions
			\end{itemize}
		\end{enumerate}
		
		\textcolor{myblue}{\textbf{Key Insight:}} Ethical committees help \textcolor{myblue}{discuss and address biases} through diverse perspectives!
	\end{frame}
	
	\begin{frame}{Slide 73: C-Suite and Board Training}
		\fontsize{6.5}{7.5}\selectfont
		\textbf{Training for Leadership:}
		
		\vspace{0.1cm}
		\textbf{Why Leadership Education Matters:}
		\begin{itemize}
			\item Wide range of risks from AI
			\item Financial, legal, reputational consequences
			\item Responsible AI is a strategic imperative
			\item Leadership drives organizational culture
		\end{itemize}
		
		\vspace{0.1cm}
		\textbf{Key Topics for Training:}
		\begin{enumerate}
			\item \textbf{AI Risks:}
			\begin{itemize}
				\item Bias and discrimination
				\item Privacy violations
				\item Regulatory non-compliance
				\item Reputational damage
			\end{itemize}
			\item \textbf{Ethical/Responsible AI Practices:}
			\begin{itemize}
				\item Fairness principles
				\item Governance structures
				\item Risk mitigation
			\end{itemize}
			\item \textbf{Strategic Considerations:}
			\begin{itemize}
				\item Competitive advantage
				\item Trust building
				\item Long-term value
			\end{itemize}
		\end{enumerate}
		
		\textcolor{myblue}{\textbf{Key Insight:}} Leadership must be \textcolor{myblue}{educated about AI risks and responsible practices}!
	\end{frame}
	
	\begin{frame}{Slide 74: Fairness and Bias Summary}
		\fontsize{6.5}{7.5}\selectfont
		\textbf{Summary: Bias, Discrimination, and Fairness}
		
		\vspace{0.1cm}
		\textbf{Key Concepts:}
		\begin{enumerate}
			\item \textbf{Bias:}
			\begin{itemize}
				\item Systemic deviation from norm
				\item Can be justifiable or problematic
				\item Sources: Problem, data, modeling, deployment
			\end{itemize}
			\item \textbf{Discrimination:}
			\begin{itemize}
				\item Bias that disfavors protected characteristics
				\item Legally prohibited
				\item Requires proactive measures
			\end{itemize}
			\item \textbf{Fairness:}
			\begin{itemize}
				\item Absence of problematic bias
				\item Group and individual fairness
				\item Cannot be fully automated
			\end{itemize}
		\end{enumerate}
		
		\textcolor{myblue}{\textbf{Key Insight:}} Bias management requires \textcolor{myblue}{continuous effort across the AI lifecycle}!
	\end{frame}
	
	\begin{frame}{Slide 75: Bias Mitigation Tools}
		\fontsize{6.5}{7.5}\selectfont
		\textbf{Tools for Bias Mitigation:}
		
		\vspace{0.1cm}
		\textbf{Available Toolkits:}
		\begin{enumerate}
			\item \textbf{Aequitas:}
			\begin{itemize}
				\item Tests models for different bias metrics
				\item Across different groups
				\item Auditing tool
			\end{itemize}
			\item \textbf{AI Fairness 360 (AIF360):}
			\begin{itemize}
				\item IBM's fairness toolkit
				\item Benchmark for testing
				\item Multiple fairness metrics
			\end{itemize}
			\item \textbf{Internal Auditing Systems:}
			\begin{itemize}
				\item Companies can create own systems
				\item Identify potential biases
				\item Continuous monitoring
			\end{itemize}
		\end{enumerate}
		
		\vspace{0.1cm}
		\textbf{External Options:}
		\begin{itemize}
			\item Third-party auditing services
			\item Independent assessors
			\item Bias reports from providers
			\item Fresh perspective
		\end{itemize}
		
		\textcolor{myblue}{\textbf{Key Insight:}} Multiple tools available for \textcolor{myblue}{identifying and mitigating bias}!
	\end{frame}
	
	% ============================================================
	% SECTION 7: PRIVACY AND CYBERSECURITY - SLIDES 76-85
	% ============================================================
	
	\begin{frame}{Slide 76: Privacy in AI Context}
		\fontsize{6.5}{7.5}\selectfont
		\textbf{Understanding Privacy in AI:}
		
		\vspace{0.1cm}
		\textbf{Definition:}
		\begin{itemize}
			\item Privacy to extent others don't have access to personal information
			\item In reference to some person or institution
			\item In reference to some personal data point
		\end{itemize}
		
		\vspace{0.1cm}
		\textbf{Why Privacy Matters:}
		\begin{enumerate}
			\item \textbf{Protection from Abuse:}
			\begin{itemize}
				\item More knowledge enables interference
				\item Power imbalances
				\item Control over individuals
			\end{itemize}
			\item \textbf{Autonomy:}
			\begin{itemize}
				\item Privacy supports self-determination
				\item Freedom from surveillance
				\item Personal space
			\end{itemize}
			\item \textbf{Dignity:}
			\begin{itemize}
				\item Respect for persons
				\item Not being reduced to data
				\item Human worth
			\end{itemize}
		\end{enumerate}
		
		\textcolor{myblue}{\textbf{Key Insight:}} Privacy = \textcolor{myblue}{lack of access to personal information}!
	\end{frame}
	
	\begin{frame}{Slide 77: Data Security}
		\fontsize{6.5}{7.5}\selectfont
		\textbf{Understanding Data Security:}
		
		\vspace{0.1cm}
		\textbf{Definition:}
		\begin{itemize}
			\item Doing what is necessary to prevent unauthorized access
			\item Protecting data from attacks and theft
			\item Ensuring data integrity
		\end{itemize}
		
		\vspace{0.1cm}
		\textbf{Key Components:}
		\begin{enumerate}
			\item \textbf{Physical Security:}
			\begin{itemize}
				\item Hardware protection
				\item Storage device security
				\item Physical access controls
			\end{itemize}
			\item \textbf{Administrative Controls:}
			\begin{itemize}
				\item Access permissions
				\item Organizational policies
				\item Employee training
			\end{itemize}
			\item \textbf{Software Applications:}
			\begin{itemize}
				\item Encryption
				\item Authentication
				\item Monitoring systems
			\end{itemize}
		\end{enumerate}
		
		\textcolor{myblue}{\textbf{Key Insight:}} Data security = \textcolor{myblue}{preventing unauthorized access} to data!
	\end{frame}
	
	\begin{frame}{Slide 78: Privacy as Ethical Issue}
		\fontsize{6.5}{7.5}\selectfont
		\textbf{Why Privacy is an Ethical Issue:}
		
		\vspace{0.1cm}
		\textbf{Wrongs:}
		\begin{itemize}
			\item Violates moral and legal rights
			\item Treats people instrumentally
			\item Reduces people to data points
			\item Respect for autonomy
		\end{itemize}
		
		\vspace{0.1cm}
		\textbf{Harms to Individuals:}
		\begin{itemize}
			\item Discrimination
			\item Blackmail
			\item Exposure and shaming
			\item Identity theft
		\end{itemize}
		
		\vspace{0.1cm}
		\textbf{Harms to Institutions:}
		\begin{itemize}
			\item Lawsuits and liability
			\item Fines and penalties
			\item Reputational damage
		\end{itemize}
		
		\vspace{0.1cm}
		\textbf{Harms to Society:}
		\begin{itemize}
			\item National security risks
			\item Democratic threats
			\item Blackmail of officials
		\end{itemize}
		
		\textcolor{myblue}{\textbf{Key Insight:}} Privacy violations cause \textcolor{myblue}{wrongs and harms} to individuals, institutions, and society!
	\end{frame}
	
	\begin{frame}{Slide 79: Personal Data as Toxic Asset}
		\fontsize{6.5}{7.5}\selectfont
		\textbf{Personal Data: A Toxic Asset}
		
		\vspace{0.1cm}
		\textbf{Why "Toxic"?}
		\begin{enumerate}
			\item \textbf{Cheap to Mine:}
			\begin{itemize}
				\item Easy to collect
				\item Low cost
				\item Widespread availability
			\end{itemize}
			\item \textbf{Valuable:}
			\begin{itemize}
				\item Profitable for companies
				\item Marketing and targeting
				\item Business insights
			\end{itemize}
			\item \textbf{Sensitive:}
			\begin{itemize}
				\item Privacy concerns
				\item Vulnerable to abuse
				\item High risk
			\end{itemize}
			\item \textbf{Hard to Keep Safe:}
			\begin{itemize}
				\item Security challenges
				\item Attacker advantage
				\item Constant threats
			\end{itemize}
		\end{enumerate}
		
		\vspace{0.1cm}
		\textbf{Key Insight:}
		\begin{itemize}
			\item Defenders are at disadvantage
			\item Attacker chooses time and method
			\item With enough resources, breach is inevitable
		\end{itemize}
		
		\textcolor{myblue}{\textbf{Key Insight:}} Personal data is a \textcolor{myblue}{toxic asset} - valuable but risky!
	\end{frame}
	
	\begin{frame}{Slide 80: Data Minimization}
		\fontsize{6.5}{7.5}\selectfont
		\textbf{The Principle of Data Minimization:}
		
		\vspace{0.1cm}
		\textbf{Core Idea:}
		\begin{itemize}
			\item Collect only necessary data
			\item Fulfill specific purpose
			\item Most effective privacy protection
		\end{itemize}
		
		\vspace{0.1cm}
		\textbf{Why It Matters:}
		\begin{enumerate}
			\item \textbf{Risk Reduction:}
			\begin{itemize}
				\item Less data = less risk
				\item Fewer targets
				\item Reduced exposure
			\end{itemize}
			\item \textbf{Respect for Persons:}
			\begin{itemize}
				\item Privacy rights
				\item Avoid unnecessary intrusion
				\item Proportionality
			\end{itemize}
			\item \textbf{Regulatory Compliance:}
			\begin{itemize}
				\item GDPR requirement
				\item Privacy laws
				\item Best practice
			\end{itemize}
		\end{enumerate}
		
		\vspace{0.1cm}
		\textbf{Key Point:}
		\begin{itemize}
			\item Only collect data when benefit outweighs disadvantages
			\item Data subjects deserve protection
			\item Minimal data collection is respectful
		\end{itemize}
		
		\textcolor{myblue}{\textbf{Key Insight:}} Data minimization = \textcolor{myblue}{collect only necessary data}!
	\end{frame}
	
	\begin{frame}{Slide 81: Right to be Forgotten}
		\fontsize{6.5}{7.5}\selectfont
		\textbf{The Right to be Forgotten:}
		
		\vspace{0.1cm}
		\textbf{Definition:}
		\begin{itemize}
			\item Person's right to have private information removed
			\item From Internet searches and other directories
			\item When data is inadequate, irrelevant, or excessive
		\end{itemize}
		
		\vspace{0.1cm}
		\textbf{When It Applies:}
		\begin{enumerate}
			\item \textbf{Incorrect Information:}
			\begin{itemize}
				\item Data that is wrong
				\item Inaccurate records
				\item False information
			\end{itemize}
			\item \textbf{Irrelevant Information:}
			\begin{itemize}
				\item No longer relevant
				\item Outdated
				\item No legitimate purpose
			\end{itemize}
			\item \textbf{Excessive Information:}
			\begin{itemize}
				\item More than needed
				\item Proportionality
				\item Over-collection
			\end{itemize}
		\end{enumerate}
		
		\textcolor{myblue}{\textbf{Key Insight:}} Right to be forgotten = \textcolor{myblue}{removal of private information when inadequate or irrelevant}!
	\end{frame}
	
	\begin{frame}{Slide 82: Control Over Data}
		\fontsize{6.5}{7.5}\selectfont
		\textbf{Giving People Control Over Their Data:}
		
		\vspace{0.1cm}
		\textbf{Key Rights:}
		\begin{enumerate}
			\item \textbf{Consent Management:}
			\begin{itemize}
				\item Easy to deny consent
				\item Easy to withdraw consent
				\item Clear options
			\end{itemize}
			\item \textbf{Access:}
			\begin{itemize}
				\item See personal data
				\item See inferences made
				\item Understand data use
			\end{itemize}
			\item \textbf{Transfer:}
			\begin{itemize}
				\item Data portability
				\item Move data to other services
				\item Freedom to choose
			\end{itemize}
			\item \textbf{Deletion:}
			\begin{itemize}
				\item Right to delete data
				\item Remove personal information
				\item Erasure
			\end{itemize}
		\end{enumerate}
		
		\textcolor{myblue}{\textbf{Key Insight:}} Control over data = \textcolor{myblue}{consent, access, transfer, and deletion rights}!
	\end{frame}
	
	\begin{frame}{Slide 83: Contextual Integrity}
		\fontsize{6.5}{7.5}\selectfont
		\textbf{Contextual Integrity in Privacy:}
		
		\vspace{0.1cm}
		\textbf{Core Idea:}
		\begin{itemize}
			\item Data use should adhere to contextual norms
			\item When data is transferred to different context, norms are violated
			\item Context matters
		\end{itemize}
		
		\vspace{0.1cm}
		\textbf{Examples:}
		\begin{enumerate}
			\item \textbf{Banking Context:}
			\begin{itemize}
				\item Data given for transaction
				\item Should not be sold to marketing company
				\item Violates expectations
			\end{itemize}
			\item \textbf{Healthcare Context:}
			\begin{itemize}
				\item Patient data given to doctor
				\item Should not end up with data brokers
				\item Violates trust
			\end{itemize}
			\item \textbf{Social Media Context:}
			\begin{itemize}
				\item Expectations about data use
				\item Context-specific norms
				\item Violations cause backlash
			\end{itemize}
		\end{enumerate}
		
		\textcolor{myblue}{\textbf{Key Insight:}} Contextual integrity = \textcolor{myblue}{data use should respect context-specific norms}!
	\end{frame}
	
	\begin{frame}{Slide 84: Data Deletion Practices}
		\fontsize{6.5}{7.5}\selectfont
		\textbf{Best Practices for Data Deletion:}
		
		\vspace{0.1cm}
		\textbf{Why Delete Data:}
		\begin{enumerate}
			\item \textbf{Protect Individuals:}
			\begin{itemize}
				\item Reduce exposure
				\item Privacy protection
				\item Security
			\end{itemize}
			\item \textbf{Maintain Accuracy:}
			\begin{itemize}
				\item Personal data changes quickly
				\item People change jobs, houses, preferences
				\item Outdated data is inaccurate
			\end{itemize}
			\item \textbf{Reduce Liability:}
			\begin{itemize}
				\item Less data = less risk
				\item Fewer targets
				\item Smaller breach impact
			\end{itemize}
		\end{enumerate}
		
		\vspace{0.1cm}
		\textbf{Best Practice:}
		\begin{itemize}
			\item Regular deletion schedules
			\item Expiry dates for data
			\item Not collected with intention of keeping forever
			\item Routine data deletion
		\end{itemize}
		
		\textcolor{myblue}{\textbf{Key Insight:}} Delete data \textcolor{myblue}{as soon as no longer necessary} - routine practice!
	\end{frame}
	
	\begin{frame}{Slide 85: Privacy Principles Summary}
		\fontsize{6.5}{7.5}\selectfont
		\textbf{Summary of Privacy Principles:}
		
		\vspace{0.1cm}
		\textbf{Key Principles:}
		\begin{enumerate}
			\item \textbf{Right to Privacy:}
			\begin{itemize}
				\item Moral and legal right
				\item Universal Declaration of Human Rights
				\item Constitutional protections
			\end{itemize}
			\item \textbf{Data Minimization:}
			\begin{itemize}
				\item Collect only necessary data
				\item Most effective protection
			\end{itemize}
			\item \textbf{Right to be Forgotten:}
			\begin{itemize}
				\item Remove inadequate data
				\item Irrelevant or excessive
			\end{itemize}
			\item \textbf{Control Over Data:}
			\begin{itemize}
				\item Consent, access, transfer, deletion
			\end{itemize}
			\item \textbf{Contextual Integrity:}
			\begin{itemize}
				\item Respect context-specific norms
			\end{itemize}
			\item \textbf{Data Deletion:}
			\begin{itemize}
				\item Delete when no longer necessary
			\end{itemize}
			\item \textbf{Data Security:}
			\begin{itemize}
				\item Encryption, anonymization, differential privacy
			\end{itemize}
		\end{enumerate}
		
		\textcolor{myblue}{\textbf{Key Insight:}} Multiple principles guide \textcolor{myblue}{ethical privacy practices}!
	\end{frame}
	
	% ============================================================
	% SECTION 8: GOVERNANCE CHALLENGES - SLIDES 86-95
	% ============================================================
	
	\begin{frame}{Slide 86: AI Governance Challenges}
		\fontsize{6.5}{7.5}\selectfont
		\textbf{Why AI is Difficult to Govern:}
		
		\vspace{0.1cm}
		\textbf{Key Challenges:}
		\begin{enumerate}
			\item \textbf{Power Asymmetries:}
			\begin{itemize}
				\item Tech companies more powerful than governments
				\item Wealth, influence, expertise
				\item Regulation difficulties
			\end{itemize}
			\item \textbf{Lack of Transparency:}
			\begin{itemize}
				\item Private company development
				\item Limited public information
				\item Safety testing unknown
			\end{itemize}
			\item \textbf{Lack of Interpretability:}
			\begin{itemize}
				\item Neural networks are "black boxes"
				\item Proxy variables obscure
				\item Hard to audit
			\end{itemize}
			\item \textbf{Lack of Ethics Structures:}
			\begin{itemize}
				\item No required ethics education
				\item No licensing
				\item No oversight committees
			\end{itemize}
		\end{enumerate}
		
		\textcolor{myblue}{\textbf{Key Insight:}} AI governance faces \textcolor{myblue}{multiple interconnected challenges}!
	\end{frame}
	
	\begin{frame}{Slide 87: Power Asymmetries}
		\fontsize{6.5}{7.5}\selectfont
		\textbf{Power Asymmetries in AI Governance:}
		
		\vspace{0.1cm}
		\textbf{The Problem:}
		\begin{itemize}
			\item Tech companies often more powerful than governments
			\item Wealth: Massive resources
			\item Influence: Lobbying and public relations
			\item Expertise: Advanced technical knowledge
		\end{itemize}
		
		\vspace{0.1cm}
		\textbf{Consequences:}
		\begin{enumerate}
			\item \textbf{Regulation Challenges:}
			\begin{itemize}
				\item Governments struggle to keep up
				\item Companies resist regulation
				\item Asymmetric negotiation
			\end{itemize}
			\item \textbf{Information Asymmetry:}
			\begin{itemize}
				\item Companies know more
				\item Governments lack technical expertise
				\item Hard to regulate effectively
			\end{itemize}
			\item \textbf{Accountability Gaps:}
			\begin{itemize}
				\item Who holds companies accountable?
				\item Regulatory capture risk
				\item Limited enforcement
			\end{itemize}
		\end{enumerate}
		
		\textcolor{myblue}{\textbf{Key Insight:}} Tech companies may be \textcolor{myblue}{more powerful than some governments}!
	\end{frame}
	
	\begin{frame}{Slide 88: Lack of Transparency}
		\fontsize{6.5}{7.5}\selectfont
		\textbf{Transparency Challenges in AI:}
		
		\vspace{0.1cm}
		\textbf{Why It's a Problem:}
		\begin{itemize}
			\item Most AI systems developed by private companies
			\item Not subject to transparency requirements
			\item Limited information about:
			\begin{itemize}
				\item Training datasets
				\item Testing and tuning
				\item Safety measures
			\end{itemize}
		\end{itemize}
		
		\vspace{0.1cm}
		\textbf{Consequences:}
		\begin{enumerate}
			\item \textbf{No Public Oversight:}
			\begin{itemize}
				\item Academics can't study systems
				\item Journalists can't investigate
				\item Public is in the dark
			\end{itemize}
			\item \textbf{Regulatory Challenges:}
			\begin{itemize}
				\item Regulators lack information
				\item Can't assess risks
				\item Enforcement difficulties
			\end{itemize}
			\item \textbf{Trust Issues:}
			\begin{itemize}
				\item Lack of transparency erodes trust
				\item Suspicion about practices
				\item Public concern
			\end{itemize}
		\end{enumerate}
		
		\textcolor{myblue}{\textbf{Key Insight:}} Private companies are \textcolor{myblue}{not subject to same transparency requirements}!
	\end{frame}
	
	\begin{frame}{Slide 89: Lack of Interpretability}
		\fontsize{6.5}{7.5}\selectfont
		\textbf{Interpretability Challenges:}
		
		\vspace{0.1cm}
		\textbf{The Black-Box Problem:}
		\begin{itemize}
			\item Neural networks are "black boxes"
			\item Even experts can't fully understand
			\item Backpropagation doesn't provide insight
			\item Hard to trace decisions
		\end{itemize}
		
		\vspace{0.1cm}
		\textbf{Why It Matters:}
		\begin{enumerate}
			\item \textbf{Accountability:}
			\begin{itemize}
				\item Can't hold accountable what you can't understand
				\item Responsibility gaps
				\item Who is responsible?
			\end{itemize}
			\item \textbf{Auditing:}
			\begin{itemize}
				\item Hard to audit black boxes
				\item New methods needed
				\item Limited effectiveness
			\end{itemize}
			\item \textbf{Governance:}
			\begin{itemize}
				\item Can't govern what you can't understand
				\item Regulatory challenges
				\item Oversight difficulties
			\end{itemize}
		\end{enumerate}
		
		\textcolor{myblue}{\textbf{Key Insight:}} Neural networks are \textcolor{myblue}{black boxes} - even experts can't fully understand them!
	\end{frame}
	
	\begin{frame}{Slide 90: Proxy Variables Challenge}
		\fontsize{6.5}{7.5}\selectfont
		\textbf{Proxy Variables and Governance:}
		
		\vspace{0.1cm}
		\textbf{What Are Proxies?}
		\begin{itemize}
			\item Techniques to simplify training
			\item Represent complex objectives with simpler metrics
			\item Example: Word count as proxy for article quality
		\end{itemize}
		
		\vspace{0.1cm}
		\textbf{The Challenge:}
		\begin{enumerate}
			\item \textbf{Alignment Issues:}
			\begin{itemize}
				\item Proxy may not align with true objective
				\item Optimization for wrong goal
				\item Unintended consequences
			\end{itemize}
			\item \textbf{Opacity:}
			\begin{itemize}
				\item If proxies are unknown
				\item Hard to understand model behavior
				\item Governance difficulties
			\end{itemize}
			\item \textbf{Scale and Complexity:}
			\begin{itemize}
				\item Billions of parameters
				\item Hard to trace processes
				\item Audit challenges
			\end{itemize}
		\end{enumerate}
		
		\textcolor{myblue}{\textbf{Key Insight:}} Unknown proxies \textcolor{myblue}{make governance difficult}!
	\end{frame}
	
	\begin{frame}{Slide 91: Lack of Ethics Structures}
		\fontsize{6.5}{7.5}\selectfont
		\textbf{Lack of AI Ethics Structures:}
		
		\vspace{0.1cm}
		\textbf{Comparison with Medical Ethics:}
		\begin{tabular}{|p{4cm}|p{4cm}|}
			\hline
			\textbf{Medical Ethics} & \textbf{AI Ethics} \\
			\hline
			Bioethicists & Few AI ethicists \\
			Ethical codes & Developing codes \\
			Ethics committees & Rare committees \\
			Required coursework & Optional courses \\
			Licensed professionals & No licensing \\
			International codes & Fragmented guidance \\
			\hline
		\end{tabular}
		
		\vspace{0.1cm}
		\textbf{Key Gap:}
		\begin{itemize}
			\item Computer scientists can graduate without ethics training
			\item No certification or licensing
			\item Boards rarely include AI ethicists
			\item No oversight committees
		\end{itemize}
		
		\textcolor{myblue}{\textbf{Key Insight:}} AI ethics lacks the \textcolor{myblue}{established structures} of medical ethics!
	\end{frame}
	
	\begin{frame}{Slide 92: Lack of Regulation}
		\fontsize{6.5}{7.5}\selectfont
		\textbf{The Regulatory Gap:}
		
		\vspace{0.1cm}
		\textbf{Current State:}
		\begin{itemize}
			\item Few national laws (China exception)
			\item No specialized agencies
			\item No government-mandated auditing
			\item Voluntary frameworks only
		\end{itemize}
		
		\vspace{0.1cm}
		\textbf{Historical Parallel:}
		\begin{itemize}
			\item FDA established for food and drugs
			\item Other regulatory agencies developed
			\item AI may follow similar path
			\item EU leading with AI Office
		\end{itemize}
		
		\vspace{0.1cm}
		\textbf{Why Regulation Matters:}
		\begin{itemize}
			\item AI is international technology
			\item Cross-border data flows
			\item Need for cooperation
			\item International standards needed
		\end{itemize}
		
		\textcolor{myblue}{\textbf{Key Insight:}} Few national AI laws exist \textcolor{myblue}{(except China)} - no specialized agencies!
	\end{frame}
	
	\begin{frame}{Slide 93: Unpredictability}
		\fontsize{6.5}{7.5}\selectfont
		\textbf{Unpredictability in AI Systems:}
		
		\vspace{0.1cm}
		\textbf{Emergent Behavior:}
		\begin{itemize}
			\item Properties not explicitly coded
			\item Capabilities that surprise
			\item Mistakes that are unexpected
			\item Interactions that are unforeseen
		\end{itemize}
		
		\vspace{0.1cm}
		\textbf{Consequences:}
		\begin{enumerate}
			\item \textbf{Safety Risks:}
			\begin{itemize}
				\item Unintended outcomes
				\item Harmful behavior
				\item System failures
			\end{itemize}
			\item \textbf{Governance Challenges:}
			\begin{itemize}
				\item Hard to regulate unpredictable systems
				\item Misuse difficult to prevent
				\item Unknown risks
			\end{itemize}
			\item \textbf{Trust Issues:}
			\begin{itemize}
				\item Surprises undermine trust
				\item Unpredictable = unreliable
				\item Adoption barriers
			\end{itemize}
		\end{enumerate}
		
		\textcolor{myblue}{\textbf{Key Insight:}} Emergent behavior = \textcolor{myblue}{unexpected capabilities not explicitly coded}!
	\end{frame}
	
	\begin{frame}{Slide 94: Truth-Tracking Limitations}
		\fontsize{6.5}{7.5}\selectfont
		\textbf{LLMs and Truth-Tracking:}
		
		\vspace{0.1cm}
		\textbf{How LLMs Work:}
		\begin{itemize}
			\item Statistical inferences and probabilistic guesses
			\item Designed for plausible responses
			\item Not based on truth or knowledge
		\end{itemize}
		
		\vspace{0.1cm}
		\textbf{Risks:}
		\begin{enumerate}
			\item \textbf{Misinformation:}
			\begin{itemize}
				\item Plausible but false information
				\item Hard to detect
				\item Spread at scale
			\end{itemize}
			\item \textbf{Physical Safety:}
			\begin{itemize}
				\item Incorrect medical advice
				\item Dangerous recommendations
				\item Real-world harm
			\end{itemize}
			\item \textbf{Legal Risks:}
			\begin{itemize}
				\item Libel and defamation
				\item Copyright issues
				\item Speech-related risks
			\end{itemize}
		\end{enumerate}
		
		\textcolor{myblue}{\textbf{Key Insight:}} LLMs give \textcolor{myblue}{plausible, not necessarily truthful} responses!
	\end{frame}
	
	\begin{frame}{Slide 95: Privacy and Copyright}
		\fontsize{6.5}{7.5}\selectfont
		\textbf{Privacy and Copyright Challenges:}
		
		\vspace{0.1cm}
		\textbf{Privacy Concerns:}
		\begin{itemize}
			\item Data scraped from internet
			\item Personal data of unsuspecting users
			\item GDPR compliance uncertain
			\item Right to access, modify, delete unclear
		\end{itemize}
		
		\vspace{0.1cm}
		\textbf{Copyright Concerns:}
		\begin{itemize}
			\item LLMs ingest copyrighted material
			\item Books, articles, images
			\item Lawsuits in process
			\item Legal uncertainty
		\end{itemize}
		
		\vspace{0.1cm}
		\textbf{Governance Nightmare:}
		\begin{itemize}
			\item Machine-created speech
			\item Company ownership
			\item Data acquisition unclear
			\item Regulatory challenges
		\end{itemize}
		
		\textcolor{myblue}{\textbf{Key Insight:}} Scraping data raises \textcolor{myblue}{privacy and copyright concerns}!
	\end{frame}
	
	% ============================================================
	% SECTION 9: REGULATORY LANDSCAPE - SLIDES 96-100
	% ============================================================
	
	\begin{frame}{Slide 96: Regulatory Overview}
		\fontsize{6.5}{7.5}\selectfont
		\textbf{Global AI Regulatory Landscape:}
		
		\vspace{0.1cm}
		\textbf{Key Players:}
		\begin{enumerate}
			\item \textbf{European Union:}
			\begin{itemize}
				\item GDPR (Privacy)
				\item DMA (Digital Markets)
				\item DSA (Digital Services)
				\item AI Act (Comprehensive AI Law)
			\end{itemize}
			\item \textbf{United States:}
			\begin{itemize}
				\item No federal privacy law
				\item State-level privacy laws
				\item AI Executive Order
				\item NIST AI RMF
			\end{itemize}
			\item \textbf{China:}
			\begin{itemize}
				\item Generative AI Measures
				\item PIPL (Privacy)
				\item Algorithmic Recommendation Rules
			\end{itemize}
		\end{enumerate}
		
		\textcolor{myblue}{\textbf{Key Insight:}} Different regions have \textcolor{myblue}{different regulatory approaches}!
	\end{frame}
	
	\begin{frame}{Slide 97: GDPR Overview}
		\fontsize{6.5}{7.5}\selectfont
		\textbf{European GDPR (General Data Protection Regulation):}
		
		\vspace{0.1cm}
		\textbf{Key Facts:}
		\begin{itemize}
			\item Implemented May 2018
			\item Regulates personal data
			\item Applies to data controllers and processors
			\item Extraterritorial jurisdiction
		\end{itemize}
		
		\vspace{0.1cm}
		\textbf{Data Subject Rights:}
		\begin{enumerate}
			\item Receive concise information
			\item Access personal data
			\item Request erasure
			\item Object to processing
		\end{enumerate}
		
		\vspace{0.1cm}
		\textbf{Obligations:}
		\begin{itemize}
			\item Notify breaches within 72 hours
			\item "Legitimate interest" required
			\item Significant fines for non-compliance
		\end{itemize}
		
		\textcolor{myblue}{\textbf{Key Insight:}} GDPR = \textcolor{myblue}{European data protection regulation} with global reach!
	\end{frame}
	
	\begin{frame}{Slide 98: EU AI Act}
		\fontsize{6.5}{7.5}\selectfont
		\textbf{The EU AI Act:}
		
		\vspace{0.1cm}
		\textbf{Key Facts:}
		\begin{itemize}
			\item Approved March 2024
			\item Council approved May 2024
			\item Establishes common regulatory framework
			\item Risk-based approach
		\end{itemize}
		
		\vspace{0.1cm}
		\textbf{Risk Categories:}
		\begin{enumerate}
			\item \textbf{Unacceptable Risk (Banned):}
			\begin{itemize}
				\item Biometric categorization using sensitive characteristics
				\item Untargeted facial image scraping
				\item Emotion recognition in workplace
				\item Social scoring
				\item Subliminal manipulation
			\end{itemize}
			\item \textbf{High-Risk:}
			\begin{itemize}
				\item Impact assessments required
				\item Thorough records
				\item Oversight measures
			\end{itemize}
			\item \textbf{Limited/Minimal Risk:}
			\begin{itemize}
				\item Transparency obligations
				\item Labeling requirements
			\end{itemize}
		\end{enumerate}
		
		\textcolor{myblue}{\textbf{Key Insight:}} EU AI Act = \textcolor{myblue}{first comprehensive AI regulation} - risk-based approach!
	\end{frame}
	
	\begin{frame}{Slide 99: US AI Regulation}
		\fontsize{6.5}{7.5}\selectfont
		\textbf{US AI Regulatory Framework:}
		
		\vspace{0.1cm}
		\textbf{Key Initiatives:}
		\begin{enumerate}
			\item \textbf{Executive Order on AI (2023):}
			\begin{itemize}
				\item Safe, secure, trustworthy development
				\item Address AI-enabled harms
				\item International cooperation
				\item Safety testing requirements
			\end{itemize}
			\item \textbf{NIST AI RMF:}
			\begin{itemize}
				\item Voluntary framework
				\item Govern, map, measure, manage
				\item Flexible and adaptable
				\item Non-prescriptive
			\end{itemize}
			\item \textbf{AI Safety and Security Board:}
			\begin{itemize}
				\item Department of Homeland Security
				\item Critical infrastructure focus
				\item Recommendations and guidance
			\end{itemize}
		\end{enumerate}
		
		\textcolor{myblue}{\textbf{Key Insight:}} US has \textcolor{myblue}{executive orders and voluntary frameworks} - no federal AI law yet!
	\end{frame}
	
	\begin{frame}{Slide 100: Responsible AI Summary}
		\fontsize{6.5}{7.5}\selectfont
		\textbf{Comprehensive Summary: Responsible AI}
		
		\vspace{0.1cm}
		\textbf{Ethical Foundations:}
		\begin{itemize}
			\item Consequentialism (outcomes and utility)
			\item Deontology (duties and rules)
			\item Virtue Ethics (character and virtues)
		\end{itemize}
		
		\vspace{0.1cm}
		\textbf{Core Principles:}
		\begin{itemize}
			\item Nonmaleficence, beneficence, justice, autonomy, explainability
		\end{itemize}
		
		\vspace{0.1cm}
		\textbf{Key Practices:}
		\begin{itemize}
			\item Bias identification and mitigation
			\item Privacy protection and data minimization
			\item Transparency and explainability
			\item Human oversight and accountability
		\end{itemize}
		
		\vspace{0.1cm}
		\textbf{Governance and Regulation:}
		\begin{itemize}
			\item Ethical committees and structures
			\item Algorithmic auditing
			\item Compliance with regulations (GDPR, EU AI Act)
			\item Continuous monitoring and improvement
		\end{itemize}
		
		\textcolor{myred}{\textbf{Exam Success:}} Understand \textcolor{myblue}{ethical frameworks, principles, governance, and regulation} for responsible AI!
	\end{frame}
	
\end{document}